\documentclass[11pt]{article}

\usepackage[utf8]{inputenc}
\usepackage[T1]{fontenc}
\usepackage{amsmath, amssymb}
\usepackage{graphicx}
\usepackage{booktabs}
\usepackage{hyperref}
\usepackage{geometry}
\usepackage{caption}
\usepackage{xcolor}
\usepackage{algorithm}
\usepackage{algorithmic}

\geometry{margin=2.5cm}

\title{Sustained Attractor Diversity in Doubt-Modulated FitzHugh--Nagumo Networks:\\
A Mean-Field Connectivity Boundary Governed by Coupling Degree}

\author{
    \textbf{Julien Chauvin} \\
    \small{\texttt{contact@cafevirtuel.org}} \\
    \small{Repository: \url{https://github.com/cafe-virtuel/Mem4ristor}}
}

% Date FIXE depuis le 2026-07-31 : \today rendait deux PDF du meme fichier
% indistinguables, donc impossible de dire quelle version un lecteur tient en main.
\date{Version 6.0.0 --- 31 July 2026}

\begin{document}

\maketitle

% ============================================================================
% ABSTRACT (~150 words)
% ============================================================================
\begin{abstract}
While coupled oscillator networks typically synchronize into homogeneous states, we present a mechanism---doubt-modulated coupling---that actively prevents synchronization and sustains attractor diversity from homogeneous initial conditions. In a modified FitzHugh--Nagumo model, a third per-node state variable, \emph{doubt} ($u$), modulates coupling \emph{polarity} through a shifted-tanh kernel $f(u) = \tanh(\pi(0.5-u)) + \delta$: a node couples attractively when its doubt is low and repulsively when it is high. Combined with a fraction $\eta$ of nodes receiving inverted stimulus (\emph{heretic} nodes) and weak noise for symmetry breaking, this sustains multi-state attractors from homogeneous initial conditions ($H_{\text{stable}} \approx 4.09 \pm 0.19$ bits on $N = 100$ lattices, 10 seeds). Our central result is an ablation: freezing the doubt dynamics collapses the network into synchrony---mean pairwise trajectory correlation jumping from $0.007$ to $0.658$, a complete separation of the two conditions with no overlap across a 30-seed replication---identifying $u$ as the primary anti-synchronization mechanism. Because this is measured on the Pearson correlation of node trajectories, it is independent of any entropy-binning choice and is the least parameter-sensitive of our findings. We then map the mechanism's \emph{limit}: a sharp connectivity boundary beyond which every local coupling correction---degree normalization included---fails to rescue diversity, which collapses into a ``dead zone.'' We show this boundary is governed by the network's \emph{coupling degree} through a mean-field sampling mechanism (harmonic degree $k_{\mathrm{harm}} = 1/\langle 1/\deg\rangle$), \emph{not} by algebraic connectivity: at fixed degree, varying $\lambda_2$ by a factor of ${\sim}27$ leaves the regime unchanged, and high-degree rings with $\lambda_2$ hundreds of times below the apparent threshold still die while equally low-$\lambda_2$ sparse rings survive. On the original Barab\'asi--Albert/Erd\H{o}s--R\'enyi dataset $\lambda_2$ appears to separate the regimes only because it covaries with degree; the empirical separation seen there ($\lambda_2 \in (2.13, 2.50)$, midpoint $2.31$) is a correlational boundary, not a causal spectral threshold. Finite-size scaling of the Binder cumulant shows the entropy loss across this boundary is a \emph{continuous crossover}, not a thermodynamic phase transition---a dichotomy a circuit-level (SPICE) implementation reproduces under analog fabrication tolerances.
\end{abstract}

% ============================================================================
% 1. INTRODUCTION
% ============================================================================
\section{Introduction}

Synchronization is a generic outcome of coupled oscillator dynamics~\cite{pikovsky2001,strogatz2000}. The Kuramoto model~\cite{acebrn2005}, distributed consensus algorithms~\cite{olfati2007}, and voter models~\cite{castellano2009} all converge toward uniform states, losing the diversity of initial conditions. While synchronization is desirable in many applications, it becomes a failure mode when heterogeneous steady states carry functional value---as in neural coding, opinion dynamics, or collective decision-making.

The problem is well characterized: on networks with homogeneous coupling, the largest Lyapunov exponent of the synchronized manifold is typically negative, making consensus a globally attracting state~\cite{jadbabaie2003}. Breaking this attraction requires either heterogeneous coupling, noise, or structural frustration~\cite{odor2024}.

We propose a mechanism based on \textbf{doubt-modulated coupling} in extended FitzHugh--Nagumo (FHN) networks. Each node carries a third state variable $u \in [0, 1]$ (the \textit{doubt} variable) that continuously modulates the sign and magnitude of inter-node coupling. When $u > 0.5$, coupling becomes repulsive, implementing polarity-modulated anti-synchronization at the single-node level---a mechanism distinct from frustrated synchronization in Kuramoto networks~\cite{odor2024}, where sign reversal is static and network-distributed rather than state-dependent and per-node. Combined with a fraction $\eta$ of nodes receiving inverted external stimulus (\textit{heretic nodes}), this architecture prevents consensus collapse from homogeneous initial conditions without requiring noise as the sole symmetry-breaking mechanism (heretics provide structural asymmetry; noise enables the Cold Start Protocol described in Section~2.3).

Our contributions are: (1)~a formal specification of the doubt-modulated FHN model and of its five architectural components (doubt-modulated coupling, heretic nodes, $1/\sqrt{N}$ scaling, adaptive doubt speed, and degree-normalized coupling); the ablation of Section~\ref{sec:ablation} measures the anti-synchronization contribution of three of them, while contribution~(4) below reports degree normalization as a negative result; (2)~the identification of the doubt variable as the \emph{primary} anti-synchronization mechanism, through an ablation measured on trajectory correlation (and thus independent of entropy binning): freezing $u$ raises the mean pairwise correlation from $0.007$ to $0.658$---a complete separation of the two conditions (Cohen's $d \approx 9$ on a 30-seed replication)---and drives the network from structured diversity into coherent chaos, while the full model sustains $H_{\text{stable}} \approx 4.09 \pm 0.19$ bits (100-bin continuous spectral entropy) on regular lattices; (3)~a sharp connectivity boundary on the mechanism's reach, which we trace to the network's \emph{coupling degree} through a mean-field sampling argument (harmonic degree $k_{\mathrm{harm}}$) rather than to algebraic connectivity: a degree-controlled experiment (Watts--Strogatz rewiring at fixed $\langle k\rangle$) and a finite-size argument (rings whose $\lambda_2 \to 0$ remain dead) establish that $\lambda_2$ is a non-causal covariate of degree, and the empirical $\lambda_2$-separation reported on the original dataset ($\lambda_2 \in (2.13, 2.50)$, midpoint $2.31$) reflects that covariation; and (4)~a negative result on power-law degree normalization: no exponent $\gamma \in [0, 1]$ in $D_{\text{eff}}(i) = D/\deg(i)^\gamma$ bridges the dead zone, consistent with a mechanism set by the \emph{number} of coupled neighbours (sampling) rather than by the per-node coupling weight.

\textbf{Roadmap.} The remainder of this paper is organized as follows. Section \ref{sec:related} positions our model within the existing literature. Section \ref{sec:model} formally defines the doubt-modulated FitzHugh--Nagumo equations, justifying the choice of hyperparameters and architectural components. Section \ref{sec:metrics} explicitly defines the evaluation metrics ($H_{\text{stable}}$, $U_4$). Section \ref{sec:results} presents our numerical results across regular and scale-free topologies, concluding with the Binder cumulant analysis of the dead zone. Finally, Section \ref{sec:discussion} discusses implications for neuromorphic design and homeostatic coupling regulation.

\subsection{Related Work}
\label{sec:related}
\textbf{Frustrated Synchronization and Chimeras.} Symmetry breaking in coupled oscillator networks is a deeply studied phenomenon. While the Kuramoto model explores synchronization through phase differences~\cite{acebrn2005}, modified models introducing phase lags or frustrated interactions have successfully generated chimera states, where coherent and incoherent domains coexist~\cite{abrams2004}. Chimera states have been characterized extensively in FitzHugh--Nagumo networks in particular: multichimera states emerge under nonlocal coupling~\cite{omelchenko2013}, persist under inhomogeneities of both the elements and the coupling topology~\cite{omelchenko2015}, and arise as noise-induced (coherence-resonance) chimeras in the excitable regime~\cite{semenova2016}; see~\cite{majhi2019} for a review of chimera states in neuronal networks. In all of these settings, however, the coupling kernel is fixed --- typically nonlocal, with quenched strength and sign. Our approach differs by making the coupling polarity \emph{state-dependent} rather than a fixed structural property: each node's instantaneous doubt level $u_i(t)$ determines the sign of its own coupling, allowing the network to dynamically rewire its effective topology.

\textbf{Adaptive and Modulated Coupling.} A growing body of work studies \emph{adaptive} dynamical networks, in which coupling weights evolve according to the nodes' own dynamics, so that structure and dynamics co-determine each other~\cite{berner2023,gross2008}. In the FitzHugh--Nagumo setting specifically, Hebbian-type plasticity rules (e.g.\ Hebb--Oja) drive transitions between travelling waves, synchrony, and chimera states by adapting the \emph{magnitude} of the links~\cite{provata2026}. Our mechanism is distinct in two respects. First, what adapts is not the magnitude but the \emph{sign} of the coupling: the doubt variable $u_i(t)$ flips node $i$ from attractive to repulsive coupling at $u_i = 0.5$. Second, the modulation is strictly \emph{local and per-node} --- driven by each node's disagreement with its own neighborhood --- rather than a pairwise learning rule on the edges. The doubt variable thus functions as a slow modulatory variable in the tradition of neuromodulation and homeostatic plasticity in computational neuroscience: it acts as a localized conflict signal, conceptually related to homeostatic regulation~\cite{friston2010}.

The closest precedent we are aware of is Zanette and Mikhailov~\cite{zanette2004}, who studied ensembles whose interactions alternate between attractive and repulsive under the control of \emph{individual internal variables}, and reported a transition between ordered and disordered states mediated by dynamical clustering. The distinction is one of causation rather than of form: there, the internal variables follow a prescribed \emph{autonomous} dynamics, independent of the state of the ensemble, so that the alternation of sign is imposed on the system from within each unit but not informed by it. Here $u_i(t)$ is driven by node $i$'s own disagreement with its neighbourhood, closing the loop between the collective state and the sign of the interaction that produces it; the frustration is earned rather than scheduled. A second, complementary line of work assigns positive and negative coupling to fixed subpopulations---conformists and contrarians~\cite{hong2011}---where the sign is a quenched property of each oscillator rather than a dynamical variable of it.

\textbf{Topology and coupling degree.} The dependence of synchronization on network topology, and in particular on the algebraic connectivity (Fiedler value, $\lambda_2$), is well established~\cite{pikovsky2001}. Our observation of a sharp regime boundary \emph{out} of functional diversity is, however, controlled not by $\lambda_2$ but by the mean coupling degree (Section~\ref{sec:lambda2}): denser local coupling drives each node's neighbourhood field toward a common mean, suppressing state heterogeneity irrespective of the global spectral structure.
% ============================================================================
% 2. MODEL
% ============================================================================
\section{Model}
\label{sec:model}

\subsection{FitzHugh--Nagumo Dynamics with Doubt}

Each node $i$ in a network of $N$ units evolves according to a modified FitzHugh--Nagumo system with a third variable $u_i$ (doubt):
\begin{align}
\frac{dv_i}{dt} &= v_i - \frac{v_i^3}{5} - w_i + I_{\text{ext},i} - \alpha_{\text{self}} \tanh(v_i) + \eta_v(t), \label{eq:dv} \\
\frac{dw_i}{dt} &= \varepsilon(v_i + a - bw_i) + dw_{\text{plasticity},i}, \label{eq:dw} \\
\frac{du_i}{dt} &= \frac{\varepsilon_{u,\text{eff}}(i)}{\tau_u}\bigl(k_u \sigma_{\text{local},i} + \sigma_{\text{baseline}} - u_i\bigr), \label{eq:du}
\end{align}
where $v_i$ is the activation potential, $w_i$ is the recovery variable, $u_i \in [0, 1]$ is the doubt variable (clamped), $\eta_v(t) \sim \mathcal{N}(0, \sigma_{\text{noise}}^2)$ is Gaussian noise, $\alpha_{\text{self}}$ is the self-inhibition gain, and $\bar{v}_{\mathcal{N}(i)} = |\mathcal{N}(i)|^{-1} \sum_{j \in \mathcal{N}(i)} v_j$ is the mean activation of neighbors of node~$i$. The local disagreement $\sigma_{\text{local},i} = |\bar{v}_{\mathcal{N}(i)} - v_i|$ measures the mismatch between a node's state and its neighborhood mean. The adaptive doubt rate is $\varepsilon_{u,\text{eff}}(i) = \varepsilon_u \cdot \min(1 + \alpha_s \cdot \sigma_{\text{local},i},\; C_{\text{cap}})$ with default parameters $\alpha_s = 2.0$, $C_{\text{cap}} = 5.0$ (see Section~\ref{sec:adaptive}); this adaptive form is active in all numerical experiments reported here. The plasticity term has two components: $dw_{\text{plasticity},i} = \lambda_{\text{learn}} \cdot \sigma_{\text{local},i} \cdot \mathbf{1}[u_i > 0.5] \cdot (1 - (w_i/w_{\text{sat}})^2) - w_i/\tau_{\text{plast}}$. The first term is activity-dependent growth gated by local disagreement and the innovation mask ($u_i > 0.5$); the second is an always-active exponential decay that drives $w$ toward zero in the absence of social drive (a novel element of this model, not present in standard FHN formulations). Note: the decay $-w_i/\tau_{\text{plast}}$ is active even when $\sigma_{\text{local}} = 0$ or $u_i \leq 0.5$.

The external input for each node depends on its type:
\begin{equation}
I_{\text{ext},i} = \begin{cases}
+I_{\text{stimulus}} + I_{\text{coupling},i} & \text{if node } i \text{ is normal,} \\
-I_{\text{stimulus}} + I_{\text{coupling},i} & \text{if node } i \text{ is a heretic node.}
\end{cases}
\label{eq:stimulus}
\end{equation}

\subsection{Doubt-Modulated Coupling}

The coupling term implements polarity-modulated anti-synchronization via a shifted-tanh kernel:
\begin{equation}
I_{\text{coupling},i} = \frac{D}{\sqrt{N}} \cdot w_i(u_i) \cdot \bigl(\bar{v}_{\mathcal{N}(i)} - v_i\bigr),
\label{eq:coupling}
\end{equation}
where the coupling weight is
\begin{equation}
w_i(u_i) = \tanh\bigl(\pi(0.5 - u_i)\bigr) + \delta, \qquad \delta = 0.01.
\label{eq:sigmoid}
\end{equation}
The factor $\pi$ ensures a slope of $-\pi$ at $u=0.5$, providing a sharp but smooth polarity transition without artificial discontinuities. For $u_i < 0.5$, $w_i > 0$ (attractive coupling: the node tends toward its neighbors). For $u_i > 0.5$, $w_i < 0$ (repulsive coupling: the node actively diverges from its neighbors). The offset $\delta = 0.01$ ensures $w_i(0.5) > 0$, preventing a hard mathematical dead zone where nodes become entirely disconnected exactly at the transition point. This small positive bias softly favors synchrony in moments of pure uncertainty. This kernel can be interpreted as a continuous analogue of frustrated interactions in spin glass systems~\cite{toulouse1977,odor2024}.

\subsection{Heretic Nodes and Structural Heterogeneity}

A fraction $\eta$ of nodes are designated as \textit{heretic nodes}: they receive their external stimulus with inverted polarity (Eq.~\ref{eq:stimulus}). This creates a permanent structural asymmetry analogous to random-field disorder in spin systems~\cite{nattermann1998}. On a $k$-regular lattice with heretic fraction $\eta$, the probability that a normal node is adjacent to at least one heretic is:
\begin{equation}
P_{\text{contact}} = 1 - (1-\eta)^k.
\label{eq:percolation}
\end{equation}
For $k=4$ (2D grid) and $\eta = 0.15$: $P_{\text{contact}} \approx 0.48$, meaning nearly half of all normal nodes experience direct heretic influence. The heuristic value $\eta = 0.15$ was empirically selected as the minimum fraction required to robustly perturb the network out of a homogeneous consensus state without overwhelmingly dominating the majority dynamics. This is consistent with the empirically observed critical threshold on regular lattices (Section~\ref{sec:results_lattice}).

\textbf{Cold Start Protocol.} Unless otherwise noted, simulations use homogeneous initial conditions ($v_i = w_i = 0$, $H(0) = 0$), including the scaling and connectivity experiments reported below. Diversity must \emph{emerge} from noise and heretic-driven symmetry breaking, not from favorable initialization.

\subsection{Adaptive Extensions}
\label{sec:adaptive}

\textbf{Adaptive doubt speed.} The fixed doubt learning rate $\varepsilon_u$ in Eq.~\ref{eq:du} can be replaced by a per-node adaptive rate modulated by local disagreement:
\begin{equation}
\varepsilon_{u,\text{eff}}(i) = \varepsilon_u \cdot \min\!\bigl(1 + \alpha_s \cdot \sigma_{\text{local}}(i),\; C_{\text{cap}}\bigr),
\label{eq:meta_doubt}
\end{equation}
with surprise gain $\alpha_s = 2.0$ and cap $C_{\text{cap}} = 5.0$. When a node's neighborhood contradicts its internal state ($\sigma_{\text{local}}$ large), its doubt updates faster. Setting $\alpha_s = 0$ recovers the fixed-rate dynamics of Eq.~\ref{eq:du}.

\textbf{Topological rewiring.} On networks with explicit adjacency matrices, nodes with $u_i > u_{\text{rewire}}$ may disconnect from their most consensual neighbor and reconnect to a random non-neighbor (degree-preserving, cooldown-regulated). This was found insufficient to resolve diversity collapse on scale-free networks alone (Section~\ref{sec:topo_regimes}).

\textbf{Degree-normalized coupling.} To address hub-induced synchronization on heterogeneous networks, the uniform coupling $D/\sqrt{N}$ is replaced with per-node weights:
\begin{equation}
D_{\text{eff}}(i) = \frac{D}{\deg(i)^\gamma}, \qquad \gamma \in [0, 1],
\label{eq:degree_norm}
\end{equation}
normalized so that $\langle D_{\text{eff}} \rangle = D/\sqrt{N}$.

\subsection{Simulation Protocol}

All simulations use Euler integration with $\Delta t = 0.05$, validated for numerical stability to $20{,}000+$ steps via systematic $\Delta t$ sweep ($\Delta t \in [0.01, 0.10]$). The stochastic term $\eta_v(t)$ is integrated with the Euler--Maruyama convention: at each step the noise increment is drawn as $\sigma_{\text{noise}} \sqrt{\Delta t}\, \xi$ with $\xi \sim \mathcal{N}(0,1)$, so that $\sigma_{\text{noise}}$ is the diffusion coefficient of the underlying stochastic differential equation rather than a per-step amplitude. All quantitative results in this manuscript were generated under this convention.

\subsection{Evaluation Metrics}
\label{sec:metrics}

To quantify the state of the network, we explicitly define two primary metrics:
\begin{enumerate}
    \item \textbf{Stable Entropy ($H_{\text{stable}}$)}: The continuous Shannon entropy of the network's activation potential distribution, computed over 100 bins (chosen to balance resolution and statistical robustness for $N=100$ networks), averaged over the last 25\% of the simulation to discard transients. It serves as our primary order parameter, bounded in $[0, \log_2(100)] \approx [0, 6.64]$ bits. A high $H_{\text{stable}}$ indicates diverse attractor states; $H_{\text{stable}} \to 0$ indicates global synchronization.
    \item \textbf{Binder Cumulant ($U_4$)}: A fourth-order statistical measure $U_4 = 1 - \frac{\langle H_{\text{stable}}^4 \rangle}{3\langle H_{\text{stable}}^2 \rangle^2}$ used in finite-size scaling to characterize the dead zone.
\end{enumerate}
A legacy 5-state discrete entropy ($H_{\text{cog}}$) using physiological $v$-thresholds (Table~\ref{tab:states}, Appendix~\ref{app:legacy_bins}) is retained only for SPICE hardware comparisons; it is not used as a primary metric in Python simulations (see Limitations, Section~\ref{sec:limitations}).

Reference parameters: $a = 0.7$, $b = 0.8$, $\varepsilon = 0.08$, $D = 0.15$, $\alpha_{\text{self}} = 0.15$, $\sigma_{\text{noise}} = 0.05$, $\sigma_{\text{baseline}} = 0.05$, $\eta = 0.15$ (see Appendix~\ref{app:parameters} for all parameters). Topologies tested: 2D periodic lattice (stencil Laplacian) and Barab\'asi--Albert ($m \in \{1, \ldots, 15\}$, $N = 100$). Watts-Strogatz and Erd\H{o}s--R\'enyi topologies were used during development for qualitative robustness checks but are not systematically reported here. Seeds and repetitions are reported per experiment.

% ============================================================================
% 3. THEORETICAL ANALYSIS
% ============================================================================
\section{Theoretical Analysis}

\subsection{Existence of Limit Cycles}

For a single \textbf{decoupled} unit ($D = 0$, $N = 1$) with $u$ treated as a quasi-static parameter ($\varepsilon_u \ll 1$), the $(v, w)$ subsystem reduces to a 2-dimensional modified FHN oscillator.

\textbf{Correction (audit 2026-04-22).} At the reference parameters $\alpha_{\text{self}} = 0.15$, $a = 0.7$, $b = 0.8$, $\varepsilon = 0.08$, the isolated unit is \emph{not} in a limit-cycle regime. Numerical fixed-point analysis locates the unique equilibrium at $v^* \approx -1.294$, $w^* \approx -0.732$ with Jacobian eigenvalues $\lambda \approx -0.055 \pm 0.283i$: a \textit{stable spiral}. The Hopf bifurcation occurs at $\alpha_{\text{crit}} \approx 0.296$; at $\alpha = 0.15$ the system is in the \textit{excitable (sub-Hopf) regime}, not oscillatory. The Poincar\'e--Bendixson argument below therefore does \emph{not} apply to the reference parameterization. It holds for $\alpha > \alpha_{\text{crit}} \approx 0.296$, and we retain it here for that parameter regime only. \textbf{Numerical verification (2026-05-05):} simulating a decoupled node ($D = 0$, $\sigma_v = 0$, $I = 0$, 3000 steps) at $\alpha = 0.15$ yields $\mathrm{std}(v) = 0.0025$ (stable spiral confirmed); at $\alpha = 0.35 > \alpha_{\text{crit}}$, $\mathrm{std}(v) = 1.22$ (sustained limit-cycle oscillation confirmed). Reproduced by \texttt{experiments/verify\_pb\_isolated\_node.py}.

By the Poincar\'e--Bendixson theorem, which applies strictly to planar ($\mathbb{R}^2$) autonomous systems (and only for $\alpha > \alpha_{\text{crit}}$):
\begin{enumerate}
\item The self-inhibition term $-\alpha_{\text{self}} \tanh(v)$ ensures boundedness of $|v|$, and $b > 0$ ensures stability of $w$. A compact positively invariant set $\Omega \subset \mathbb{R}^2$ exists.
\item For $\alpha > \alpha_{\text{crit}} \approx 0.296$, the unique equilibrium is unstable (Jacobian eigenvalues have positive real parts).
\item By Poincar\'e--Bendixson, the absence of a stable equilibrium within $\Omega$ guarantees the existence of a limit cycle for the \emph{isolated unit}.
\end{enumerate}

\textbf{Important scope note.} The Poincar\'e--Bendixson theorem does \emph{not} extend to the coupled $3N$-dimensional system, and we make no analytical stability claim for the full network. Linearization around the synchronized manifold at the reference parameterization in fact yields Jacobian eigenvalues with \emph{negative} real parts ($\approx -0.13$), i.e.\ perfect synchrony would be locally \emph{stable} in a homogeneous network without heretics (Section~\ref{sec:comp_analysis}). The observed sustained diversity is therefore \emph{not} derived from an analytical instability of the synchronized state: it is an empirical result, produced structurally by heretic asymmetry and symmetry-breaking initialization (Section~\ref{sec:comp_analysis}) and quantified by the ablation of Section~\ref{sec:ablation}.


\subsection{Heretic Percolation Threshold}

The heretic fraction $\eta = 0.15$ is an \textbf{empirical observation}, not a derived constant. Its effectiveness on regular lattices admits a qualitative interpretation via Eq.~\ref{eq:percolation}: at $\eta = 0.15$ on a 4-regular graph, 48\% of normal nodes are adjacent to at least one heretic, providing sufficient local frustration to prevent consensus nucleation.

An analogy with the random-field Ising model~\cite{nattermann1998} suggests that the critical heretic concentration should scale inversely with mean degree: $\eta_c \approx C / \langle k \rangle$. Fitting $C$ to match the observed threshold on 2D lattices ($\langle k \rangle = 4$, $\eta_c \approx 0.15$) gives $C \approx 0.6$. This is an empirical fit, not a derivation, and should be treated as a guideline. In particular, on scale-free networks, the degree heterogeneity invalidates the mean-field assumption underlying this estimate.

\subsection{Entropy Lower Bound}

We sketch a proof that $H > 0$ whenever $\eta > 0$ \textbf{and $I_{\text{stimulus}} \neq 0$}. Suppose the system reaches consensus ($v_i = v^*$ for all $i$). Then the Laplacian coupling vanishes. Normal nodes receive $+I_{\text{stimulus}}$ while heretics receive $-I_{\text{stimulus}}$, driving them toward different equilibria. This contradicts the consensus assumption, so consensus is not an attractor when $\eta > 0$ and $I_{\text{stimulus}} \neq 0$. A quantitative lower bound via mean-field analysis yields:
\begin{equation}
H \geq -\eta \log_2 \eta - (1 - \eta) \log_2(1 - \eta) \approx 0.61 \text{ bits for } \eta = 0.15,
\label{eq:entropy_bound}
\end{equation}
representing the mixing entropy of the two subpopulations.

\textbf{Scope note (audit 2026-04-22).} The argument above is \emph{vacuous} when $I_{\text{stimulus}} = 0$: at this baseline, inverting polarity is a no-op because both normal and heretic nodes receive zero direct stimulus drive. This zero-stimulus setting is used exclusively for the \emph{endogenous} connectivity and dead-zone experiments (Sections~\ref{sec:topo_regimes}--\ref{sec:lambda2}), where the (limited) diversity arises from coupling-mediated asymmetry and noise rather than from heretic inversion. The sustained-diversity results lie outside this regime: Table~\ref{tab:scaling} ($H_{\text{stable}} \approx 4.09$ bits; synchrony $= 0.036$) is evaluated under the driven protocol ($I_{\text{stimulus}} = 0.5$), where heretics are active and the lower bound of Eq.~\ref{eq:entropy_bound} applies.

% ============================================================================
% 4. RESULTS
% ============================================================================
\section{Results}
\label{sec:results}

The central result of this section is causal: doubt dynamics are what keep the network out of synchrony. Section~\ref{sec:results_lattice} first establishes that sustained diversity exists on regular lattices; Section~\ref{sec:ablation} then shows, by ablation, that removing the doubt variable collapses this diversity into synchrony---a complete separation of the two conditions in pairwise trajectory correlation ($0.007$ to $0.658$, Cohen's $d \approx 9$), measured independently of any entropy-binning choice, making it our least parameter-sensitive finding. The remaining subsections map where this mechanism reaches its limit: a connectivity boundary, governed by coupling degree, beyond which no local coupling correction rescues diversity.

\subsection{Diversity on Regular Lattices}
\label{sec:results_lattice}

Table~\ref{tab:scaling} reports the sustained attractor entropy across network scales on 2D periodic lattices. All runs use the Cold Start Protocol ($v = w = 0$), $I_{\text{stimulus}} = 0.5$, $\eta = 0.15$, 3000 steps, 10 seeds $\in \{42, 123, 777, 17, 256, 1337, 99, 314, 2024, 888\}$. Results reproduced by \texttt{experiments/verify\_table1\_preprint.py}.

\begin{table}[ht]
\centering
\caption{Scaling performance on 2D periodic lattices ($I_\mathrm{stim}=0.5$, $\eta = 0.15$, 10 seeds $\in\{42,123,777,17,256,1337,99,314,2024,888\}$, 3000 steps, last 25\% window). $H_{\text{stable}}$: 100-bin continuous spectral entropy. Mean doubt $\bar{u}$: mean of the per-node doubt variable; $\bar{u} \approx 1$ indicates fully active repulsive coupling, consistent with state-driven chimera-like behavior~\cite{abrams2004} (see Section~\ref{sec:model}); note that the mechanism here is state-driven via $u_i(t)$, mechanistically distinct from the quenched non-local coupling of the original chimera formulation~\cite{abrams2004}. All values reproduced by \texttt{experiments/verify\_table1\_preprint.py}.}
\label{tab:scaling}
\begin{tabular}{lccc}
\toprule
\textbf{Metric} & \textbf{$4 \times 4$} & \textbf{$10 \times 10$} & \textbf{$25 \times 25$} \\
\midrule
$N$                          & 16   & 100   & 625   \\
$H_{\text{stable}}$ (bits)   & $3.21 \pm 0.12$ & $4.09 \pm 0.19$  & $4.40 \pm 0.09$  \\
Mean doubt $\bar{u}$          & 0.996 & 0.988 & 0.984 \\
\bottomrule
\end{tabular}
\end{table}

The system sustains substantial voltage diversity ($H_{\text{stable}} > 3.2$ bits) across all tested scales, with entropy increasing with network size as more independent oscillatory modes emerge. Note: the initial 200--500 time steps exhibit transient entropy peaks before converging to the sustained regime. All $H_{\text{stable}}$ values in this paper use 100-bin continuous spectral entropy unless explicitly noted.

\subsection{Component Ablation}
\label{sec:ablation}

Table~\ref{tab:ablations} quantifies the contribution of each mechanism to diversity preservation. This ablation is the paper's central causal result: it isolates the doubt dynamics as the primary anti-synchronization engine, on a metric (pairwise Pearson correlation of trajectories) that does not depend on any entropy-binning convention.

\begin{table}[ht]
\centering
\caption{Ablation study (BA $m=3$, $N=100$, \texttt{degree\_linear}, 3000 steps, 10 seeds, $I_\mathrm{stim}=0.5$). Primary metric: pairwise synchrony (mean Pearson correlation of $v(t)$ trajectories; lower = more independent = more diverse). LZ: mean temporal Lempel--Ziv complexity of node state sequences. Snapshot entropy ($H_{\mathrm{cog}}$, $H_{\mathrm{cont}}$) is not reported here as it measures voltage spread at a single instant, not behavioral independence, and gives directionally incorrect results for this comparison.}
\label{tab:ablations}
\begin{tabular}{lccc}
\toprule
\textbf{Configuration} & \textbf{Synchrony} $\downarrow$ & \textbf{LZ complexity} & \textbf{Collapse?} \\
\midrule
Full model                     & $0.002 \pm 0.009$ & $1.096 \pm 0.022$ & No \\
No heretics ($\eta = 0$)       & $0.034 \pm 0.023$ & $1.154 \pm 0.026$ & Partial \\
No Levitating Sigmoid          & $0.003 \pm 0.008$ & $1.069 \pm 0.030$ & No \\
Frozen $u$ (no doubt dynamics) & $0.697 \pm 0.085$ & $1.602 \pm 0.005$ & Yes \\
\bottomrule
\end{tabular}
\end{table}

Removing the heretic mechanism raises synchrony ($0.002 \to 0.034$): nodes begin to co-evolve. LZ complexity also rises ($1.096 \to 1.154$), meaning the trajectories become \emph{less} predictable rather than more structured (Section~\ref{sec:traj_minimality}: $C_{\text{LZ}} \approx 0$ is structured, $\approx 1$ is a random walk). The two metrics therefore do not move in the same direction here, and only synchrony supports the loss of independence. Removing the Levitating Sigmoid, by contrast, leaves synchrony unchanged ($0.002 \to 0.003$, within seed-to-seed variation) and slightly lowers LZ complexity: on this metric and in this regime, it is not individually necessary. Freezing $u$ entirely produces a catastrophic synchronization surge into the fully locked regime. Because the Full-model baseline is statistically indistinguishable from zero (its standard deviation exceeds its mean), a fold-change ratio is unstable; we therefore report the effect as a \emph{complete separation} of the two synchrony distributions (Cohen's $d \approx 9$ on a 30-seed replication, \texttt{experiments/b4\_ablation\_robustness.py}), confirming that the doubt dynamics are the primary anti-synchronization engine.

\subsection{Trajectory-Based Minimality: Coordination Phase Space}
\label{sec:traj_minimality}

Snapshot entropy ($H_{100}$, $H_{\text{cog5}}$) measures spatial diversity at a single instant. It conflates \emph{structured diversity} (nodes follow distinct but predictable trajectories) with \emph{random disorder} (nodes explore independently and chaotically). To resolve this ambiguity, we introduce two complementary trajectory metrics computed on the full $v(t)$ history:
\begin{itemize}
  \item \textbf{Pairwise synchrony} $\bar{r}$: mean Pearson cross-correlation of $v(t)$ traces between all node pairs (subsampled). $\bar{r} \approx +1$: globally locked; $\bar{r} \approx 0$: independent; $\bar{r} \approx -1$: anti-phase.
  \item \textbf{LZ76 complexity} $C_{\text{LZ}}$: normalised Lempel--Ziv complexity of per-node state sequences ($\approx 0$: structured/predictable; $\approx 1$: random walk).
\end{itemize}

Figure~\ref{fig:phase_space} plots all 80 runs of the ablation study (4 conditions $\times$ 2 protocols $\times$ 10 seeds) in the $(\bar{r}, C_{\text{LZ}})$ plane. The four quadrants define interpretable dynamical regimes.

\begin{figure}[ht]
\centering
\includegraphics[width=0.82\textwidth]{../../../figures/coordination_phase_space.png}
\caption{Coordination phase space $(\bar{r}, C_{\text{LZ}})$ for four ablation conditions under two protocols (BA $m=3$, $N=100$, \texttt{degree\_linear}, 10 seeds each). Each point is one run; filled markers are centroids $\pm$ std. Quadrant boundaries: $\bar{r}=0.15$, $C_{\text{LZ}}=1.6$. \textbf{FULL} is the only condition that occupies the structured half-plane ($C_{\text{LZ}} < 1.6$) in both protocols. FROZEN\_U migrates to the chaotic quadrant (high sync, high LZ) under forcing---a signature of \emph{coherent chaos}. Ablating $u$ (FROZEN\_U) produces strongly seed-dependent trajectories with high LZ regardless of protocol.}
\label{fig:phase_space}
\end{figure}

The key claim is that \textbf{FROZEN\_U is the only condition lying clearly outside the structured ($C_{\text{LZ}} < 1.6$) half-plane in both protocols} (0/10 seeds below the boundary in each). NO\_HERETIC leaves it under forcing only ($1.77$). The full model and NO\_SIGMOID sit inside it under forcing ($1.55$ and $1.54$) and \emph{at its boundary} under the endogenous protocol ($1.62$ and $1.60$): the full model is therefore at the edge of the structured half-plane rather than distinctively within it, consistent with the near-criticality reported in Section~\ref{sec:discussion}. Note also that under the endogenous protocol the FULL and NO\_HERETIC conditions are identical by construction---at $I_{\text{stim}} = 0$ the heretic inversion is a no-op (Section~\ref{sec:model})---so that protocol contrasts three distinct systems, not four. This is still a stronger minimality criterion than snapshot entropy: the phase-space position captures the temporal organisation of dynamics, not merely their instantaneous diversity.

\textbf{Complementary evidence from forcing sweep.} To demonstrate that $u$ acts specifically as an \emph{anti-synchronisation filter}, we swept $I_{\text{stim}} \in [0, 1]$ for FULL vs FROZEN\_U (BA $m=3$, $N=100$, 7 seeds, 3000 steps). FULL maintains $\bar{r} \approx 0.025$--$0.035$ across all $I_{\text{stim}} \in [0.1, 1.0]$ (flat, near-zero synchrony). FROZEN\_U shows a sharp transition at $I_{\text{stim}} \approx 0.2$ from $\bar{r} = 0.006$ (endogenous) to $\bar{r} = 0.830$ (plateau), with Cohen's $d = 13.21$ ($p = 9.7 \times 10^{-10}$) separating the two conditions at $I_{\text{stim}} = 0.5$. This confirms that $u$ absorbs the common forcing signal and prevents stimulus-induced phase locking---a \emph{dynamic buffering} function that snapshot metrics cannot detect.

\subsection{Doubt as an Active Decorrelator: Spatial Mutual Information}
\label{sec:mi_analysis}

To characterize the role of $u$ beyond snapshot entropy, we computed the pairwise mutual information MI$(d)$ as a function of graph hop-distance $d$ under four ablation conditions on Lattice ($N=100$) and BA $m=3$ ($N=100$). MI is estimated via 2D joint histograms on $v(t)$ histories post-warmup ($I_{\mathrm{stim}}=0.5$, \texttt{degree\_linear}, 10 seeds).

\begin{center}
\begin{tabular}{lcccc}
\toprule
Ablation & MI($d=1$) lattice & Decay & MI($d=1$) BA & Decay \\
\midrule
FULL       & 0.601 & 0.092 & 0.488 & $-0.193$ \\
NO\_HERETIC & 1.003 & 0.091 & 0.577 & $-0.124$ \\
FROZEN\_U  & \textbf{1.712} & 0.057 & \textbf{1.662} & 0.306 \\
NO\_SIGMOID & 0.577 & 0.091 & 0.459 & $-0.174$ \\
\bottomrule
\end{tabular}
\end{center}

The key finding: \textbf{freezing $u$ (FROZEN\_U) sharply raises mutual information}, by $2.85\times$ on the 2D lattice (0.601 $\to$ 1.712 nats) and by $3.41\times$ on BA $m=3$ (0.488 $\to$ 1.662 nats; 10 seeds in both protocols), indicating that the $u$ dynamics actively decorrelate neighboring nodes. The FULL and NO\_SIGMOID conditions yield comparably low MI, so the decorrelation is attributable to the $u$ dynamics themselves rather than to the sigmoid coupling kernel alone. On BA $m=3$, the negative decay coefficients indicate that MI under decorrelated conditions \emph{increases} mildly with hop distance, a hub-mediated effect absent on the regular lattice. This complements the trajectory result of Section~\ref{sec:traj_minimality}: $u$ functions as a dynamic buffering mechanism both in the time domain (preventing phase-locking under forcing) and in the spatial domain (reducing pairwise MI between neighbors).

\subsection{Topological Normalization Regimes}
\label{sec:topo_regimes}

We systematically sweep the Barab\'asi--Albert attachment parameter $m \in \{1, 2, 3, 4, 5, 6, 7, 8, 10, 15\}$ with $N = 100$, 10 seeds per configuration, 3000 steps, and $I_{\text{stimulus}} = 0$. Two coupling modes are compared: uniform ($D/\sqrt{N}$, i.e.\ $\gamma = 0$) and degree-linear ($D/\deg(i)$, i.e.\ $\gamma = 1$). Note: Table~\ref{tab:alpha_sweep} extends this analysis to intermediate $\gamma$ values.

\begin{table}[ht]
\centering
\caption{Regime map (Cold Start Protocol, endogenous $I_{\text{stimulus}}=0$, 10 seeds): BA attachment parameter $m$ vs.\ coupling normalization. $H_{\text{cog}}^{\text{uni/dl}}$ (5-bin, Appendix~\ref{app:legacy_bins}) is a \emph{relative} multi-state indicator, uniform vs.\ degree-linear; absolute values are not cited. The functional co-metrics for degree-linear coupling---continuous entropy $H_{\text{cont}}^{\text{dl}}$ and pairwise synchrony $\bar{r}^{\text{dl}}$---are shown alongside (full uniform columns in the CSV). The regime boundary is governed by the coupling degree (Section~\ref{sec:lambda2}); the apparent $\lambda_2 \approx 2.31$ separation on the initial hand-labelled dataset is a correlational proxy, not a causal spectral threshold. Consult the measured-regime regression (Figure~\ref{fig:regime_degree}) for the primary, binning-independent evidence. Source: \texttt{figures/limit02\_regime\_map\_functional.csv}.}
\label{tab:ba_m_sweep}
\begin{tabular}{rrrrrrl}
\toprule
$m$ & $\lambda_2$ & $H_{\text{cog}}^{\text{uni}}$ & $H_{\text{cog}}^{\text{dl}}$ & $H_{\text{cont}}^{\text{dl}}$ & $\bar{r}^{\text{dl}}$ & Regime \\
\midrule
1  & 0.02  & 0.67 & 0.80 & 3.09 & $-0.00$ & both functional \\
2  & 0.63  & 0.27 & 0.71 & 3.14 & $0.00$  & degree-linear \\
3  & 1.40  & 0.01 & 0.68 & 3.14 & $0.01$  & degree-linear \\
4  & 2.23  & 0.00 & 0.62 & 3.20 & $0.03$  & degree-linear \\
5  & 3.08  & 0.00 & 0.32 & 3.33 & $0.13$  & degree-linear (weak) \\
6  & 4.01  & 0.00 & 0.03 & 2.85 & $0.10$  & dead-zone onset \\
7  & 4.91  & 0.00 & 0.00 & 2.60 & $-0.00$ & dead zone \\
8  & 5.81  & 0.00 & 0.00 & 2.46 & $-0.00$ & dead zone \\
10 & 7.78  & 0.00 & 0.00 & 2.31 & $-0.00$ & dead zone \\
15 & 12.58 & 0.00 & 0.00 & 2.09 & $-0.00$ & dead zone \\
\bottomrule
\end{tabular}
\end{table}

Two \textbf{regime transitions} are observed in the multi-state indicator:
\begin{enumerate}
\item $m = 1 \to m = 2$: uniform coupling loses functional diversity (tree topology) while degree-linear becomes and stays dominant through the sparse regime.
\item $m = 5 \to m = 6$: a gradual entry into a dead zone where \emph{neither normalization sustains multi-state diversity}. Under the Cold Start Protocol this onset sits near $m \approx 6$, later than the $m \approx 5$ reported by earlier non-cold-start runs.
\end{enumerate}

In this Barab\'asi--Albert family the spectral gap $\lambda_2$ tracks the transition closely, but only because it grows monotonically with the attachment parameter $m$ (hence with mean degree); Section~\ref{sec:lambda2} shows by degree-controlled experiments that the causal predictor is the coupling degree, not $\lambda_2$. Note that pairwise synchrony $\bar{r}$ remains near zero across the entire sweep (peak $0.13$): the endogenous dead zone is a spatial collapse onto a common fixed point, not a temporal consensus---synchrony discriminates regimes only under external drive (Section~\ref{sec:ablation}).

\subsection{Power-Law Exponent Sweep}
\label{sec:alpha_sweep}

To test whether an intermediate exponent $\gamma$ in $D_{\text{eff}}(i) = D/\deg(i)^\gamma$ (Eq.~\ref{eq:degree_norm}) can bridge the dead zone identified in Table~\ref{tab:ba_m_sweep}, we sweep $\gamma \in [0, 1]$ across $m \in \{2, 3, 4, 5, 6, 8, 10\}$ (10 seeds, 3000 steps, $I_{\text{stimulus}} = 0$). Note that this sweep probes the \emph{endogenous} regime ($I_{\text{stimulus}} = 0$, heretic inversion inactive), a stricter condition than the driven protocol ($I_{\text{stimulus}} = 0.5$) of the main results.

\begin{table}[ht]
\centering
\caption{Sweep of the normalization exponent $\gamma$, endogenous regime ($I_{\text{stimulus}}=0$), Cold Start Protocol, 10 seeds, $\gamma \in [0,1]$ in steps of $0.1$. For each $m$ we report the exponent $\gamma^*$ maximizing the multi-state indicator $H_{\text{cog}}$ (5-bin, Appendix~\ref{app:legacy_bins}), with the functional co-metrics at that $\gamma^*$: continuous spectral entropy $H_{\text{cont}}$ and mean pairwise synchrony $\bar{r}$. No exponent sustains multi-state diversity beyond $m \approx 5$, and the decline is gradual; synchrony stays near zero throughout, i.e.\ the endogenous regime never locks into temporal consensus (see metric note below). Source: \texttt{figures/limit02\_alpha\_sweep.csv}.}
\label{tab:alpha_sweep}
\begin{tabular}{rrrrrl}
\toprule
$m$ & $\gamma^*$ & $\max H_{\text{cog}}$ & $H_{\text{cont}}$ & $\bar{r}$ & Assessment \\
\midrule
2  & 0.6 & 0.72 & 3.16 & $-0.00$ & Multi-state diversity \\
3  & 0.9 & 0.69 & 3.13 & $0.01$  & Multi-state diversity \\
4  & 1.0 & 0.62 & 3.20 & $0.03$  & Multi-state diversity \\
5  & 1.0 & 0.32 & 3.33 & $0.13$  & Weak / declining \\
6  & 1.0 & 0.03 & 2.85 & $0.10$  & Dead-zone onset \\
8  & --- & $\approx 0$ & 2.49 & $-0.00$ & Dead zone \\
10 & --- & $\approx 0$ & 2.31 & $-0.00$ & Dead zone \\
\bottomrule
\end{tabular}
\end{table}

The negative result holds under the continuous metric: \textbf{no exponent $\gamma$ recovers multi-state diversity at high degree}. At the best $\gamma$ for each $m$, $H_{\text{cog}}$ declines gradually (0.72 at $m=2$ to $\approx 0$ by $m \gtrsim 6$)---a smooth crossover, not the abrupt $m=2 \to 3$ collapse reported by earlier non-cold-start runs, which the Cold Start Protocol (Section~\ref{sec:results_lattice}) resolves. The continuous entropy $H_{\text{cont}}$ declines even more gently (from $\approx 3.2$ to $\approx 2.3$ bits), confirming that the loss is one of \emph{multi-state} structure rather than of all voltage spread. Crucially, synchrony $\bar{r}$ stays near zero for every $m$ and $\gamma$ (peak $0.13$, versus $0.697$ for the frozen-doubt ablation of Section~\ref{sec:ablation}): the endogenous dead zone is \emph{not} a temporal consensus but a spatial collapse onto a common fixed point. Under the driven protocol ($I_{\text{stimulus}} = 0.5$, heretics active), functional diversity extends further; the endogenous regime is more fragile, consistent with the heretic mechanism being one of the two jointly necessary ingredients identified in Section~\ref{sec:comp_analysis}.

\textbf{Metric note (why $H_{\text{cog}}$ here).} We re-measured this sweep with the paper's functional metrics to check whether the artifact-prone $H_{\text{cog}}$ could be retired. It cannot, for a specific reason: in the endogenous regime no continuous metric cleanly reproduces the multi-state boundary. Pairwise synchrony---the primary metric of the ablation study (Section~\ref{sec:ablation})---is uninformative here because, absent a common drive, nodes never synchronize regardless of degree ($\bar{r} \approx 0$ everywhere). Continuous entropy $H_{\text{cont}}$ is misleading at weak coupling: near-decoupled nodes (low $\gamma$, low $m$) spread widely in voltage and inflate $H_{\text{cont}}$ without forming distinct functional states---at $m=2$, $H_{\text{cont}}$ peaks at $\gamma=0$ while $H_{\text{cog}}$ peaks at $\gamma=1$. The multi-state boundary is therefore intrinsically discrete, and $H_{\text{cog}}$ is retained as a \emph{relative} indicator of it, with absolute values not cited. The robust, binning-independent regime evidence is carried elsewhere: the measured-regime regression of Figure~\ref{fig:regime_degree} (continuous entropy across topologies at comparable coupling) and the synchrony-based ablation of Section~\ref{sec:ablation} under drive.

Why no exponent suffices is developed next: degree-based normalization ($D/\deg^\gamma$) reweights each coupling signal but cannot change the \emph{number} of neighbours averaged into a node's mean-field target---the sampling quantity that actually sets the regime (Section~\ref{sec:lambda2}). Reweighting cannot undo a boundary fixed by how many samples enter the local average.

\subsection{Locating the Boundary: Coupling Degree, not Algebraic Connectivity}
\label{sec:lambda2}

The regime boundary of Section~\ref{sec:topo_regimes} coincides, on the Barab\'asi--Albert and Erd\H{o}s--R\'enyi families, with a threshold in algebraic connectivity near $\lambda_2 \approx 2.3$. Labelling each of 36 configurations from 12 topology types ($N \approx 100$) as functional ($H_{\mathrm{cog}} > 0$) or dead ($H_{\mathrm{cog}} \approx 0$) from the normalization sweeps of Section~\ref{sec:topo_regimes} yields an apparent \emph{complete separation} in $\lambda_2$: functional configurations at $\lambda_2 < 2.13$, dead configurations at $\lambda_2 > 2.50$, midpoint $2.31$ (source: \texttt{figures/p2\_edge\_betweenness.csv}). Two facts forbid reading this as a spectral threshold. First, the separation is \emph{correlational}: the regime labels are assigned by topology type, and $\lambda_2$ covaries with mean degree across the sampled families, so a clean $\lambda_2$-split is expected whether or not $\lambda_2$ is causal. Second, the gap is thin and its upper edge rests on a single atypical seed, so the ``$2.31$'' point estimate carries a wide, unreported uncertainty. We therefore treat $\lambda_2 \approx 2.31$ as an empirical, correlational boundary on this dataset and test its causality directly.

\textbf{A degree-controlled test rules out $\lambda_2$ as the cause.} To break the degree--$\lambda_2$ covariation we ran two controlled experiments (scripts in \texttt{experiments/lambda2\_foundation\_20260701/}). (i)~\emph{Fixed degree, swept $\lambda_2$}: Watts--Strogatz rewiring holds the mean degree $\langle k\rangle$ exactly constant while the rewiring probability sweeps $\lambda_2$ over a ${\sim}27\times$ range. At each fixed degree the dead/functional outcome is \emph{flat} in $\lambda_2$; a $k=10$ ring ($\lambda_2 = 0.22 \ll 2.31$) dies in 100\% of runs, exactly like a random-regular graph of the same degree at $\lambda_2 = 4.5$. (ii)~\emph{Finite-size invariance}: for a $k=8$ ring, $\lambda_2$ falls from $0.118$ to $0.0074$ as $N$ grows $100 \to 400$ (roughly $300\times$ below the apparent threshold), yet the dead fraction stays at $1.0$ throughout, while a $k=4$ ring stays functional over the same $\lambda_2$ range. A quantity that vanishes with $N$ cannot cause a size-independent effect. Both results are incompatible with a causal role for $\lambda_2$.

\textbf{Mechanism: mean-field sampling.} The coupling term pulls each node toward the mean activation of its neighbours---an average of $\deg(i)$ samples of the local field---so the dispersion of these targets across nodes scales as $\langle 1/\deg\rangle$ (law of large numbers). A low-degree node samples few neighbours and keeps a noisy, distinct target (diversity survives); a high-degree node averages many neighbours, its target collapses onto the global mean, and the network locks into consensus (cognitive death). A kinematic test (targets vs.\ predictors, no dynamics) gives correlation $r = 0.9998$ between the target variance and $\langle 1/\deg\rangle$, against $r = -0.50$ for $\lambda_2$.

\textbf{Regression on measured regimes.} To test this on the dynamics rather than the kinematics, and without the hand-assigned labels of the initial dataset, we simulated 70 networks (14 topologies $\times$ 5 seeds: Barab\'asi--Albert $m \in \{2,\dots,6\}$, Erd\H{o}s--R\'enyi, random-regular and ring families; $N = 200$, endogenous regime, cold start) and labelled each \emph{by measurement}, using the continuous 100-bin spectral entropy $H_{\mathrm{cont}}$---not the 5-bin $H_{\mathrm{cog}}$, whose absolute values are treated as unreliable throughout this work---rather than by topology type (Figure~\ref{fig:regime_degree}). Regressing the measured regime on each candidate predictor, the coupling degree tracks $H_{\mathrm{cont}}$ far better than $\lambda_2$: Spearman $\rho = -0.83$ for the harmonic degree $k_{\mathrm{harm}} = 1/\langle 1/\deg\rangle$ and $-0.78$ for the mean degree, against only $-0.52$ for $\lambda_2$ (Pearson $r = -0.80$, $-0.75$, $-0.57$ respectively; all $p < 10^{-6}$, $n = 70$). The four families collapse onto a single monotonic curve versus $k_{\mathrm{harm}}$, whereas $\lambda_2$ scatters them---most starkly the rings, whose $\lambda_2 \approx 0$ spans the full range of outcomes as their degree grows. A binary dead/functional split labelled by the legacy $H_{\mathrm{cog}}$ reproduces the same \emph{ordering} (misclassifying $4/70$ configurations for $k_{\mathrm{harm}}$, $8/70$ for mean degree, and $17/70$ for $\lambda_2$), consistent with the earlier hand-labelled separation.

Two caveats temper the numerical boundary. First, under the continuous metric the regime is a \emph{smooth, monotonic decline} (from $\approx 3.9$ to $\approx 2.6$ bits across the sweep), not a collapse to zero: the sharp ``dead zone'' is in part an artifact of the discrete 5-bin metric, and no clean binary threshold exists in $H_{\mathrm{cont}}$ (only $5/70$ instances fall below the largest gap). Second, $k_{\mathrm{harm}}$ and mean degree are close under $H_{\mathrm{cont}}$ ($-0.83$ vs.\ $-0.78$): the data robustly identify the \emph{coupling degree}, not $\lambda_2$, as the driver, but do not sharply single out the harmonic degree over the plain mean degree (that distinction is clearer under $H_{\mathrm{cog}}$, and rests on the mechanism rather than on this fit alone). The picture nonetheless reconciles degree heterogeneity: $k_{\mathrm{harm}}$ is dominated by low-degree nodes (Jensen: $\langle 1/\deg\rangle \ge 1/\langle k\rangle$), so scale-free networks retain noisy targets through their low-degree periphery and survive to higher mean degree than a regular graph.

\begin{figure}[htbp]
  \centering
  \includegraphics[width=\linewidth]{../../../figures/a3_regime_regression_hcont.png}
  \caption{\textbf{Regime measured on real simulations, regressed on degree vs.\ algebraic connectivity.} Each point is one simulated network (14 topologies $\times$ 5 seeds, $N = 200$, endogenous regime, cold start), its regime measured as the continuous 100-bin spectral entropy $H_{\mathrm{cont}}$---not assigned by topology type. (\textit{Left}) Versus the harmonic degree $k_{\mathrm{harm}} = 1/\langle 1/\deg\rangle$, the four families (Barab\'asi--Albert, Erd\H{o}s--R\'enyi, random-regular, ring) collapse onto a single monotonic curve (Spearman $\rho = -0.83$). (\textit{Right}) Versus algebraic connectivity $\lambda_2$, the same points scatter ($\rho = -0.52$); the rings (red), all near $\lambda_2 \approx 0$, span the full range of outcomes as their degree grows, showing $\lambda_2$ carries no causal signal. Dashed line: the $H_{\mathrm{cont}}$ level below which only $5/70$ instances fall---there is no clean binary boundary, the decline is continuous. Source: \texttt{experiments/a3\_regime\_regression\_hcont.py}.}
  \label{fig:regime_degree}
\end{figure}

The dead zone is thus a real and robust property of the network, but its origin is the coupling degree, not the spectral structure: denser coupling suppresses diversity by increasing the number of neighbours each node averages. The numerical boundary ($k_{\mathrm{harm}} \approx 6$) depends on the regime metric and is reported as an empirical value; the causal claim---degree, not $\lambda_2$---does not.

\subsection{Adaptive Coupling Protocol $D(u) = D_{\max} \cdot u$}
\label{sec:adaptive_coupling}

The degree-normalization failure at high coupling degree raises the question whether a protocol \emph{different in kind}, not merely in weighting, could extend the functional regime. We test an adaptive coupling rule in which the coupling strength itself scales with the local doubt variable:

\begin{equation}
D_i(u_i) = D_{\max} \cdot u_i, \qquad D_{\max} = 0.50.
\label{eq:adaptive_D}
\end{equation}

This creates a positive feedback loop: rising doubt amplifies coupling strength, which increases local disagreement, which further raises doubt. The feedback is stabilizing for sparse topologies (doubt saturates at intermediate values) and destabilizing for dense topologies (doubt saturates at $u \approx 1$, collapsing consensus).

\textbf{Method.} Barabási--Albert networks $m \in \{1, 2, 3, 4, 5, 6, 7, 8\}$, $N = 100$, 5 seeds, $T = 10{,}000$, \texttt{degree\_linear} normalization ($\gamma = 1$), $I_{\text{stimulus}} = 0$. Three protocols compared: $D = 0$ (no coupling), $D = 0.15$ (reference constant), $D(u) = 0.50 \cdot u$ (adaptive).

\begin{table}[ht]
\centering
\caption{Protocol comparison: $D(u) = 0.50 \cdot u$ vs.\ $D = 0.15$ across BA attachment parameter $m$. $H_{\text{cont}}$: 100-bin continuous spectral entropy. $\bar{r}$: mean synchrony. Best result per row in bold.}
\label{tab:adaptive_D_sweep}
\begin{tabular}{rrrrrrr}
\toprule
 & \multicolumn{2}{c}{$D(u) = 0.50 \cdot u$} & \multicolumn{2}{c}{$D = 0.15$} & \multicolumn{2}{c}{$D = 0$} \\
$m$ & $H_{\text{cont}}$ & $\bar{r}$ & $H_{\text{cont}}$ & $\bar{r}$ & $H_{\text{cont}}$ & $\bar{r}$ \\
\midrule
1 & \textbf{3.70} & \textbf{0.72} & 3.68 & 0.51 & 3.34 & 0.71 \\
2 & \textbf{4.36} & 0.39 & 4.71 & 0.02 & 3.34 & 0.71 \\
3 & 2.85 & 0.01 & \textbf{4.19} & $-$0.01 & 3.34 & 0.71 \\
4 & 1.55 & $-$0.00 & \textbf{3.90} & $-$0.01 & 3.34 & 0.71 \\
5 & 0.18 & $-$0.01 & \textbf{3.32} & $-$0.01 & 3.34 & 0.71 \\
6 & 0.00 & $-$0.01 & \textbf{2.42} & $-$0.01 & 3.34 & 0.71 \\
7 & 0.00 & $-$0.01 & 1.31 & $-$0.01 & 3.34 & 0.71 \\
8 & 0.00 & $-$0.01 & 0.67 & $-$0.01 & 3.34 & 0.71 \\
\bottomrule
\end{tabular}
\end{table}

\textbf{Results} (Table~\ref{tab:adaptive_D_sweep}):

\begin{enumerate}
\item $D = 0$ is inoperative across all $m$ --- dead zone is absolute without adaptive coupling. $\bar{r} \approx 0.71$ constant, $H_{\text{cont}} \approx 3.34$ (noise-only regime).
\item $D = 0.15$ provides the most \emph{consistent} performance: gradual decline from $m = 2$ onward, functional up to $m = 6$.
\item $D(u) = 0.50 \cdot u$ achieves the \emph{best single result} at $m = 1$ ($\bar{r} = 0.72$, $H = 3.70$) and dominates at $m = 2$ ($H = 4.36$). However, it exhibits \textbf{catastrophic collapse at $m \geq 5$}: $H \to 0$, $\bar{r} \to -0.01$ (full consensus lock-in). Mean doubt $\bar{u} \approx 1.0$ at these points, confirming saturation of the feedback loop.
\end{enumerate}

\textbf{Interpretation.} The adaptive protocol exploits the same feedback mechanism in opposite directions depending on topology. In sparse networks ($m = 1$--$2$), doubt rises just enough to amplify anti-synchronization without triggering saturation. In dense networks ($m \geq 5$), the denser signaling pathways drive $u$ to its upper bound faster, locking the system into a state where coupling is permanently maximal and repulsive --- the network freezes into forced consensus. The collapse is not an artifact of the entropy metric: synchrony drops to $\bar{r} = -0.01$ alongside $H \to 0$, indicating genuine \textbf{consensus closure by doubt saturation}.

\textbf{Protocol recommendation.} The optimal coupling strategy is topology-dependent:
\begin{itemize}
\item $m \leq 5$ (sparse to intermediate): use $D(u) = 0.50 \cdot u$ for best diversity and synchrony decoupling;
\item $m \geq 6$ (dense scale-free): fall back to $D = 0.15$ constant, which remains functional through the dead zone;
\item $D = 0$ is never functional.
\end{itemize}

Critically, the regime boundary $m \approx 5$ is \textbf{not a universal topological invariant} --- it depends on the chosen coupling protocol. This nuance refines the claim in Section~\ref{sec:topo_regimes}: the dead zone for degree-normalized coupling ($m \geq 5$) is a property of that specific protocol family, not of the model alone.

\subsection{Finite-Size Scaling of the Dead Zone}
\label{sec:binder_cumulant}

The regime boundary identified in Section~\ref{sec:topo_regimes} raises a natural question: is the entry into the dead zone a genuine thermodynamic phase transition, or a smooth crossover? To answer it, we apply the standard order-of-transition diagnostic, the Binder cumulant~\cite{binder1981finite},

\begin{equation}
U_4(\lambda_2) = 1 - \frac{\langle H_{\text{stable}}^4 \rangle}{3\langle H_{\text{stable}}^2 \rangle^2},
\end{equation}

where $H_{\text{stable}}$ is the stable entropy defined in Section~\ref{sec:metrics} and $\langle \cdot \rangle$ denotes the ensemble average over independent initial conditions and topological realizations at fixed algebraic connectivity $\lambda_2$. For each $\lambda_2$ value, the network is evolved for $T_{\text{eq}} = 1875$ time units to discard transients, and statistics are collected over $T_{\text{stat}} = 625$ time units. In the finite-size-scaling (FSS) framework, a crossing or converging minimum of $U_4(\lambda_2)$ across system sizes signals a critical point; a featureless $U_4$ indicates the absence of one.

\textbf{Campaigns.} Three successive campaigns of increasing scope inform this analysis. An initial pilot ($N_{\text{real}} = 9$ across 4 $\lambda_2$ bins; \texttt{figures/v6\_binder\_cumulant\_U4.csv}) was too small to support FSS conclusions and is retained only for provenance. Campaign~J ($N_{\text{real}} = 1800$ simulations, $N \in \{100, 200, 400, 800, 1600\}$, BA $m = 3$; \texttt{figures/campaign\_j\_agg.csv}) provides the size-resolved panel of Figure~\ref{fig:binder_cumulant}. Finally, an extended cartography ($N_{\text{real}} = 701$ valid runs, 6 seeds per configuration; \texttt{figures/v6\_binder\_dryrun\_extended\_U4.csv}) spans $m \in \{1, 2, 5, 7\}$, $N \in \{100, 200\}$, and $\lambda_2 \in [0.5, 10.0]$ (Figure~\ref{fig:binder_cartography}), covering sparse, intermediate, and dense topologies on both sides of the regime boundary.

\textbf{Result: no thermodynamic transition.} Across both independent campaigns, $U_4$ remains essentially flat at the Gaussian reference value $2/3$. In the extended cartography, the amplitude of $U_4$ variations never exceeds $0.012$ --- an order of magnitude below the $0.05$--$0.20$ range expected for a clear Binder crossing~\cite{binder1981finite} --- with no detectable sign change or crossing of the $2/3$ value, uniformly across sparse ($m = 1$, $\lambda_2 \in [0.5, 2.5]$), intermediate ($m = 2$, $\lambda_2 \in [0.6, 4.6]$), and dense ($m = 5, 7$, $\lambda_2 \in [2.9, 10.0]$) regimes. In Campaign~J, no converging minimum is observed across system sizes; the deepest (still small) $U_4$ deviations at large $N$ occur at $\lambda_2 \approx 7$--$8$, far from the regime boundary. The order parameter $\langle H_{\text{stable}} \rangle$, in contrast, decreases smoothly and monotonically across the same region (a drop of $0.5$--$2.0$ bits depending on $m$, steepest near $\lambda_2 \approx 2$ for $m \in \{1, 2\}$). Error bars in both figures denote one standard deviation across realizations; no bootstrap resampling is applied, as the ensemble size is treated as a fixed computational budget.

\textbf{Interpretation.} The combination of a flat $U_4$ and a smooth, monotonic $\langle H_{\text{stable}} \rangle$ is the diagnostic signature of a \emph{crossover}, not a phase transition in the thermodynamic sense. We therefore abandon the $U_4$-based order-of-transition diagnostic for this system and adopt $\langle H_{\text{stable}} \rangle$ as the primary order parameter (see also Section~\ref{sec:limitations}). This reading is fully compatible with the sharp \emph{classification} boundary of Section~\ref{sec:lambda2}: the empirical separation near $\lambda_2 \approx 2.31$ cleanly sorts functional from dead-zone behavior on the sampled families, but---as shown there---this is a correlational proxy for the coupling degree, and the entropy decrease through the boundary is continuous rather than critical.

\begin{figure}[ht]
\centering
\includegraphics[width=\textwidth]{../../../figures/v6_binder_cumulant.png}
\caption{
\textbf{Finite-size scaling of the dead zone.}
(\textit{Left}) Binder cumulant $U_4 = 1 - \langle H_{\text{stable}}^4 \rangle / (3\langle H_{\text{stable}}^2 \rangle^2)$ as a function of the algebraic connectivity $\lambda_2$ for network sizes $N \in \{100, 200, 400, 800, 1600\}$.
Each data point averages independent realizations after transient discard.
With increasing $N$, the $U_4$ minimum broadens and flattens, approaching a plateau characteristic of a continuous cross-over, centered near $\lambda_2 \approx 2.31$ (dashed red line).
(\textit{Right}) Corresponding order parameter $\langle H_{\text{stable}} \rangle$ showing a size-independent entropy drop across the crossover region, signaling the progressive loss of functional degrees of freedom in the dense regime.
The inset displays the numerical derivative $|d\langle H \rangle / d\lambda_2|$ for $N=1600$, showing a smooth peak near $\lambda_2 \approx 2.3$ rather than a critical divergence.
}
\label{fig:binder_cumulant}
\end{figure}

\begin{figure}[ht]
\centering
\includegraphics[width=\textwidth]{../../../figures/v6_binder_dryrun_extended.png}
\caption{
\textbf{Extended cartography of the Binder cumulant across topologies ($m \in \{1, 2, 5, 7\}$, $N \in \{100, 200\}$, $N_{\text{real}} = 701$).}
(\textit{Top row}) Binder cumulant $U_4$ as a function of the algebraic connectivity $\lambda_2$ for each attachment parameter $m$. The dashed red line marks the canonical $\lambda_{2,\text{crit}} \approx 2.31$, the dotted gray line marks the Gaussian reference $U_4 = 2/3$. Across all four topologies and the full $\lambda_2 \in [0.5, 10.0]$ range, $U_4$ remains within $[0.653, 0.667]$ (amplitude $< 0.012$), with no detectable crossing. This rules out a first-order or two-step Binder transition over the topology range explored.
(\textit{Bottom row}) Corresponding order parameter $\langle H_{\text{stable}} \rangle$. The entropy decreases monotonically with $\lambda_2$, with the steepest decline near $\lambda_2 \approx 2$ for $m \in \{1, 2\}$ and a milder decrease for $m \in \{5, 7\}$. The combination of a flat $U_4$ and a monotonic $\langle H_{\text{stable}} \rangle$ is the diagnostic signature of a smooth cross-over, not a thermodynamic phase transition.
}
\label{fig:binder_cartography}
\end{figure}

\subsection{LZ76 Regime Classification: Beyond the Binder Cumulant}
\label{sec:lz_regime}

While the Binder cumulant $U_4$ (§
\ref{sec:binder_cumulant}) does not exhibit a convergent minimum, the Lempel--Ziv complexity $C_{\text{LZ}}$ of node trajectories provides a complementary, structurally meaningful classification of the network state. We systematically swept the Barab\'asi--Albert attachment parameter $m \in \{3, 4, 5, 6, 7, 8, 9, 10\}$ at $N = 100$ with $5$ seeds per configuration, comparing three coupling protocols: (i) $D = 0$ (uncoupled baseline), (ii) $D = 0.15$ (static reference), and (iii) $D(u) = 0.50 \cdot u$ (adaptive, this work). All runs use $\texttt{degree\_linear}$ normalization, $I_{\text{stim}} = 0$, and $T = 2000$ steps.

\begin{figure}[ht]
\centering
\includegraphics[width=\textwidth]{../../../figures/fss_lz_regime_map.png}
\caption{
\textbf{LZ76 regime classification across coupling protocols.}
Three-panel heatmap of the Lempel--Ziv complexity $C_{\text{LZ}}$ as a function of Barab\'asi--Albert attachment $m$ (rows) and Fiedler value $\lambda_2$ (columns), for three coupling protocols (panels): (a) $D = 0$ (uncoupled baseline, all $C_{\text{LZ}} \approx 1.10$ --- pure FHN chaos); (b) $D = 0.15$ static (gradual decrease from $C_{\text{LZ}} \approx 0.88$ at $m=3$ to $C_{\text{LZ}} \approx 0.78$ at $m=10$); (c) $D(u) = 0.50 \cdot u$ adaptive (sharp transition at $m=6$ from $C_{\text{LZ}} = 0.88$ to $C_{\text{LZ}} = 0.65$, $-27\%$, followed by stable plateau at $C_{\text{LZ}} \approx 0.58$ through $m = 10$).
The adaptive protocol uniquely produces a \emph{structured dead zone}: trajectories remain temporally organized ($C_{\text{LZ}} < 0.85$) even as the degree-driven crossover in mean entropy $H_{\text{stable}}$ reduces overall diversity.
The LZ-based classification thus identifies a regime boundary that the Binder cumulant cannot resolve.
Reproduced by \texttt{experiments/fss\_lz\_sweep.py} (150 simulations, 5 seeds, ~85~s).
}
\label{fig:lz_regime}
\end{figure}

The key finding, summarized in Table~
\ref{tab:lz_regime}: the adaptive protocol $D(u) = 0.50 \cdot u$ undergoes a sharp transition at $m = 6$ (LZ drops from $0.88$ to $0.65$, $-27\%$), entering a regime where temporal trajectories are significantly more structured ($C_{\text{LZ}} < 0.85$) despite the degree-driven crossover in mean entropy. This structured regime persists up to $m = 10$ ($C_{\text{LZ}} = 0.58$), whereas the static protocol $D = 0.15$ exhibits $C_{\text{LZ}} > 0.83$ throughout and degrades monotonically.

\begin{table}[ht]
\centering
\caption{LZ76 regime classification (BA attachment $m$, $N=100$, $5$ seeds, $I_{\mathrm{stim}}=0$, $T=2000$). $\langle C_{\mathrm{LZ}} \rangle$ and $\langle H_{\mathrm{stable}} \rangle$ are seed-averaged. The structured regime ($C_{\mathrm{LZ}} < 0.85$) is reached at $m = 6$ for the adaptive protocol but never for the static protocol in the swept range.}
\label{tab:lz_regime}
\begin{tabular}{lcccc}
\toprule
\textbf{Protocol} & $m$ & $\lambda_2$ & $\langle H_{\mathrm{stable}} \rangle$ & $\langle C_{\mathrm{LZ}} \rangle$ \\
\midrule
$D=0.50\cdot u$ (adaptive) & 5 & $3.01$ & $3.51$ & $0.88$ \\
$D=0.50\cdot u$ (adaptive) & \textbf{6} & $3.87$ & $\mathbf{4.23}$ & $\mathbf{0.65}$ \\
$D=0.50\cdot u$ (adaptive) & 7 & $4.83$ & $4.00$ & $0.63$ \\
$D=0.50\cdot u$ (adaptive) & 10 & $7.59$ & $3.06$ & $0.58$ \\
\midrule
$D=0.15$ (static)          & 5 & $3.01$ & $3.28$ & $0.86$ \\
$D=0.15$ (static)          & 6 & $3.96$ & $3.01$ & $0.87$ \\
$D=0.15$ (static)          & 10 & $7.70$ & $2.58$ & $0.78$ \\
\midrule
$D=0$ (uncoupled)          & 5 & $3.01$ & $4.75$ & $1.11$ \\
$D=0$ (uncoupled)          & 10 & $7.59$ & $4.73$ & $1.10$ \\
\bottomrule
\end{tabular}
\end{table}

This LZ-based regime classification complements the crossover analysis: while the mean entropy $H_{\mathrm{stable}}$ progressively decreases with $m$ for both static and adaptive protocols, only the adaptive protocol exhibits a bifurcation in temporal structure at $m = 6$. The threshold $C_{\mathrm{LZ}} < 0.85$ thus serves as a complementary, structurally meaningful regime boundary that the Binder cumulant cannot resolve. The result is reproducible: the same pattern was obtained with $3$ seeds (initial run, 2026-05-30) and $5$ seeds (re-run, 2026-06-01) to within $0.01$ on all reported metrics.

\subsection{Comparative Benchmarks}

Table~\ref{tab:benchmarks} compares the model against classical collective dynamics models on $10 \times 10$ lattices. The model's own values are those of Table~\ref{tab:scaling} (3000 steps, $I_{\text{stim}} = 0.5$, 10 seeds); the baseline entries are the documented limiting behaviours of those models rather than measurements made here (see caption). The comparison should be read as a positioning statement: these baselines converge \emph{by design}, and the point of the row is that a mechanism which prevents convergence is not a variant of them.

\begin{table}[ht]
\centering
\caption{Benchmark comparison on a $10\times10$ periodic lattice ($N = 100$, $I_\mathrm{stim} = 0.5$, 3000 steps, 10 seeds --- the same runs as Table~\ref{tab:scaling}, so that both columns of this row come from one protocol: both are written by a single execution of \texttt{experiments/verify\_table1\_preprint.py} under the Cold Start Protocol). Synchrony: mean pairwise Pearson correlation of $v(t)$ trajectories (lower = more independent). $H_{\text{stable}}$: 100-bin continuous spectral entropy. On BA $m=3$ under the ablation protocol the model reaches an even lower synchrony ($0.002 \pm 0.009$, Table~\ref{tab:ablations}), the lattice being the more conservative choice here. \textbf{What the baseline rows are, and are not.} The three baseline entries are \emph{documented limiting behaviours} of those models---phase locking above the critical coupling~\cite{acebrn2005}, the absorbing state of the voter model~\cite{castellano2009}, and convergence by construction in distributed consensus~\cite{olfati2007}---not measurements performed in this work. \textbf{No entropy is reported for them (n/a), deliberately:} $H_{\text{stable}}$ is computed on the voltage distribution, whereas Kuramoto and voter dynamics carry phases and discrete opinions, so a common entropy scale does not exist. An independent lattice-matched run (\texttt{experiments/benchmark\_kuramoto.py}) illustrates the difficulty rather than resolving it: under a shared 5-bin binning it returns $H = 2.22$ for Kuramoto against $1.73$ for this model, but it bins \emph{phases} in one case and \emph{voltages} in the other, so the ordering it produces is an artefact of the binning in either direction. The row that carries the comparison is therefore synchrony, which is the same quantity---the mean pairwise correlation of state trajectories---for every model listed.}
\label{tab:benchmarks}
\begin{tabular}{lccc}
\toprule
\textbf{Model} & \textbf{Synchrony} $\downarrow$ & $H_{\text{stable}}$ & \textbf{Diversity sustained?} \\
\midrule
Doubt-modulated FHN (this work) & $0.036 \pm 0.023$ & $4.09 \pm 0.19$ & Yes \\
Kuramoto ($K = 2$)              & $\approx 1.0$     & n/a             & No (phase-locked) \\
Voter model                     & $\approx 1.0$     & n/a             & No (absorbing state) \\
Distributed consensus           & $\approx 1.0$     & n/a             & No (by design) \\
\bottomrule
\end{tabular}
\end{table}

\subsection{SPICE Circuit Validation}
\label{sec:spice_validation}

To validate the crossover on a circuit-level substrate, we translated the FHN network model into NGSpice netlists using behavioral voltage sources for the nonlinear dynamics, and ran transient simulations on Barab\'{a}si--Albert graphs ($m = 5$, $N = 64$, degree-linear coupling). Three operating conditions were tested over 50 independent random seeds each:

\begin{table}[ht]
\centering
\caption{SPICE transient simulation results (BA $m=5$, $N=64$, degree-linear coupling, 50 seeds per condition). $H_{\text{cont}}$: 100-bin continuous spectral entropy of terminal voltage distribution; $H_{\text{cog}}$: 5-state discrete entropy (KIMI bins) retained for circuit-scale reference.}
\label{tab:spice_validation}
\begin{tabular}{lccc}
\toprule
\textbf{Condition} & $\eta$ & $\sigma_C$ & $H_{\text{cont}}$ (mean $\pm$ std) \\
\midrule
A: Dead zone (high degree, $m=5$) & 0.10 & 0.00 & $1.38 \pm 0.04$ \\
B: Functional (noise only)                              & 0.50 & 0.00 & $4.30 \pm 0.19$ \\
C: Functional (noise + CMOS mismatch)                  & 0.50 & 0.10 & $4.33 \pm 0.17$ \\
\bottomrule
\end{tabular}
\end{table}

The dead zone condition yields $H_{\text{cont}} = 1.38 \pm 0.04$ bits ($N_{\text{seeds}}=50$), consistent with collapsed diversity. The functional conditions yield $H_{\text{cont}} \approx 4.3$ bits ($3.1\times$ ratio), confirming the crossover on a circuit-level substrate. Critically, CMOS mismatch ($\sigma_C = 0.10$) does not degrade diversity ($4.33 \pm 0.17$ vs $4.30 \pm 0.19$ bits, condition B vs C), indicating that analog fabrication tolerances are compatible with the mechanism. Simulations were executed in batch mode (\texttt{ngspice -b}) on NGSpice 46, with raw results archived in \texttt{experiments/spice/results/}.

% ============================================================================
% 5. COMPUTATIONAL ANALYSIS
% ============================================================================
\section{Computational Analysis}
\label{sec:comp_analysis}

\subsection{Complexity}

On stencil-based lattices, each time step requires $\approx 56N$ floating-point operations (5-point Laplacian: $9N$; sigmoid: $4N$; FHN dynamics and doubt: $43N$), yielding $O(N)$ per-step complexity. Empirical benchmarks confirm:
\begin{equation}
t_{\text{step}}(\text{ms}) \approx 2.87 \times 10^{-4} N + 0.212,
\label{eq:scaling}
\end{equation}
enabling real-time simulation up to $N \approx 10{,}000$ on commodity hardware. For explicit adjacency matrices, Laplacian computation is $O(N^2)$, but incremental updates reduce per-rewire cost to $O(1)$. For $N > 1000$, the adjacency matrix is automatically converted to compressed sparse row (CSR) format, yielding $455\times$ memory reduction at $N = 5000$.

\subsection{Numerical Stability}

The model uses first-order Euler integration with $\Delta t = 0.05$. A systematic sweep ($\Delta t \in \{0.01, 0.02, 0.05, 0.10\}$, 20{,}000 steps) confirmed that $\Delta t \leq 0.05$ is numerically stable: after an initial transient convergence ($\sim$500 steps), entropy standard deviation drops to $\sim$0.016 and remains flat. Higher-order methods (Runge--Kutta 4) improve precision but are not required for stability.

\subsection{Origin of Sustained Diversity}

The sustained attractor diversity observed in Section~\ref{sec:results_lattice} is structural in origin rather than arising from analytical instability of a synchronized state. At the reference parameterization ($\alpha = 0.15$, sub-Hopf regime), linearization around the synchronized manifold yields Jacobian eigenvalues with negative real parts ($\approx -0.13$), indicating that perfect synchrony would be locally stable in a homogeneous network without heretics. Sustained diversity is instead produced by two mechanisms acting in concert: (1) heretic nodes receiving inverted stimulus ($-I_{\text{stim}}$) provide permanent structural asymmetry that prevents consensus nucleation; (2) the Cold Start Protocol (Gaussian noise at $t=0$, $\sigma_v = 0.1$) breaks the symmetry of homogeneous initial conditions. These two mechanisms — structural heterogeneity and symmetry-breaking initialization — are jointly necessary and sufficient to produce the chimera-like attractors~\cite{abrams2004} reported in Table~\ref{tab:scaling}. The $u$ dynamics then sustain diversity by actively decorrelating neighbours (Section~\ref{sec:mi_analysis}), as confirmed by the $3.41\times$ MI increase under FROZEN\_U ablation on BA $m=3$ ($2.85\times$ on 2D lattice).

% ============================================================================
% 6. DISCUSSION
% ============================================================================
\section{Discussion}
\label{sec:discussion}

\subsection{Polarity-Modulated Anti-Synchronization}

The doubt-modulated coupling mechanism implements polarity-modulated anti-synchronization at the single-node level: when a node's doubt exceeds $u = 0.5$, its coupling flips from attractive to repulsive. Unlike frustrated synchronization in Kuramoto-type models~\cite{odor2024}---where frustration arises from quenched random coupling signs distributed statically across the network---the polarity reversal here is a continuous, state-dependent process driven by each node's internal dynamics. The analogy to spin-glass frustration~\cite{toulouse1977} therefore holds only in its qualitative consequence (prevention of ground-state ordering), not in its mechanism. The heretic fraction provides a complementary source of structural asymmetry through permanent stimulus inversion, analogous to quenched disorder in random-field models~\cite{nattermann1998}.

\subsection{The Connectivity Boundary and Coupling Degree}

The most significant finding is the sharp regime boundary at $m \approx 5$. Section~\ref{sec:lambda2} identifies its origin: it is set by the coupling degree, not by algebraic connectivity. The mechanism is a mean-field averaging effect---each node relaxes toward the mean of $\deg(i)$ neighbours, and above a critical degree that mean collapses onto the global consensus regardless of how the coupling weight is normalized. This explains the negative result of Section~\ref{sec:alpha_sweep}: no per-node weighting $D/\deg(i)^\gamma$ bridges the dead zone, because reweighting a sum of many samples does not restore the variance lost to averaging---only reducing the \emph{number} of samples (the effective degree) or injecting node-specific drive can.

This reframes the design implication. Resolving diversity collapse on dense networks is not a matter of smarter global normalization (eigenvector centrality, Laplacian whitening): any scheme that preserves each node's full neighbourhood average inherits the same mean-field collapse. The levers that remain are structural---sparsifying the effective coupling neighbourhood---or dynamical---per-node stimulus heterogeneity that survives averaging. Both remain open for future work.

\textbf{Doubt time-scale and criticality.} A systematic sweep of $\tau_u \in [0.05, 100]$ on BA $m=3$ (N=100, \texttt{degree\_linear}, $I_{\mathrm{stim}}=0.5$, 3 seeds) reveals a sharp bifurcation. Note the direction of the parameter: $\tau_u$ enters Eq.~\ref{eq:du} as a divisor, so a \emph{larger} $\tau_u$ makes the doubt relaxation \emph{slower}. For $\tau_u < 10$ (fast doubt), frustration freezes and $H_{\mathrm{cog}} \approx 0$ (pseudo-consensus). For $\tau_u \in [10, 50]$, the system enters a transition regime ("breathing chimera"): $H_{\mathrm{cog}}$ rises from $0.023$ to $0.380$. For $\tau_u > 50$ (slow doubt), the system regains coordinated diversity ($H_{\mathrm{cog}} \approx 0.8$, $\mathrm{sync} \approx 0$). We report the regimes as measured and do not attach a mechanism to either end: an earlier reading of this sweep had the two ends inverted. \textbf{The reference value $\tau_u = 10$ sits precisely at the bifurcation edge}, which explains the bimodal distribution of outcomes observed across seeds (Section~\ref{sec:topo_regimes}). This also clarifies why the trajectory-metrics study (Section~\ref{sec:traj_minimality}) found the FULL model at the boundary of the structured half-plane: it is operating near criticality.

\subsection{Thermodynamic Viability and Homeostatic Coupling Regulation}

A fundamental question regarding the observed dead zone is whether the collapse of network dynamics represents a mere computational artifact or a physical reality of the coupled system. Our results suggest the latter. The $u_i(t)$ parameter (constitutional doubt) can be interpreted as a homeostatic regulator: each node continuously adjusts its coupling strength to reduce local disagreement, implementing a form of conflict-driven modulation without a generative model.

In this homeostatic interpretation, self-organizing systems maintain their structural and functional integrity by continually reducing local disagreement. In the Mem4ristor network, the absolute value of the local Laplacian $|\mathcal{L}_v|_i$ acts as a local conflict signal, which drives the adaptation of coupling strength for node $i$ relative to its topological neighborhood.

The dynamics of $u_i$ (Eq.~\ref{eq:du}) are explicitly driven by this conflict signal:
\begin{equation*}
\tau_u \dot{u}_i = \epsilon_u(t) \left( k_u |\mathcal{L}_v|_i + \sigma_{\text{baseline}} - u_i \right)
\end{equation*}

When a node experiences high conflict with its neighbors (high $|\mathcal{L}_v|_i$), the conflict signal increases, driving $u_i \to 1$. This increase in $u_i$ locally reduces the coupling strength $D_{\text{eff}}$, effectively decoupling the node from the source of the conflict. This dynamic implements homeostatic coupling regulation: the node alters its effective coupling with the network topology to reduce local disagreement, without invoking a generative model or variational inference.

This homeostatic mechanism is not imposed externally, but emerges from the self-regulation of the network. Preliminary numerical and circuit-level (SPICE) validation of a Kirchhoff-based self-regulation mechanism (ART; scripts and netlists available in the project repository) suggests that the effective coupling each node experiences need not be fixed but can be locally down-regulated by pruning the weights of conflicted edges, keeping the node's effective sampling neighbourhood below the critical degree in the functional regime. A full characterization of this mechanism is beyond the scope of the present paper.

The homeostatic mechanism also has physical limits: active decoupling is only viable within a bounded region of the state space. When the coupling degree overwhelms the dissipative capacity of the local $u_i$-mediated decoupling (i.e., when the effective sampling degree exceeds its critical value, Section~\ref{sec:lambda2}), local down-regulation can no longer keep each node's neighbourhood mean distinct from the global consensus. The conflict signal $|\mathcal{L}_v|$ becomes uniformly high across the network, forcing $u_i \to 1$ globally.

We term this global saturation a \textbf{consensus overload}: the network enters a state of reduced functional diversity where the continuous entropy drops toward a residual baseline ($H \to H_{\mathrm{res}} \approx 2.4$), consistent with progressive entropy depression rather than a sharp phase transition. Crucially, the Binder cumulant analysis in Section~\ref{sec:binder_cumulant} does \emph{not} provide evidence of a critical point at the connectivity boundary: no converging minimum is observed across system sizes, and the deepest $U_4$ deviations occur at $\lambda_2 \approx 7$--$8$, not at $\lambda_2 \approx 2.31$. The entropy depression is therefore best characterized as a continuous, degree-driven crossover, not a thermodynamic phase transition.

Thus, the transition into the dead zone is not a failure of the model, but a demonstration of a physical limit. The homeostatic interpretation provides an intuitive account of a phenomenon whose necessity follows from the mean-field averaging argument of Section~\ref{sec:lambda2}. The conclusion is physical: a high-degree network stripped of its capacity for localized decoupling enters a state of reduced diversity, consistent with the degree-driven crossover characterization.



\subsection{Scope and Applicability}

The doubt-modulated FHN model was originally motivated by deliberative scenarios in which premature consensus is undesirable (e.g., citizens' assemblies, collective decision-making). While the model captures the dynamical requirements for diversity preservation, it has not been validated against real-world deliberative data. The results presented here are purely computational and should be understood as properties of the dynamical system, not as claims about social behavior.

\section{Conclusion}

We have presented a FitzHugh--Nagumo network model in which doubt-modulated coupling---a per-node variable that flips coupling polarity from attractive to repulsive---together with structural heretic nodes sustains attractor diversity from homogeneous initial conditions. On regular lattices ($10\times10$, $I_{\text{stim}}=0.5$) the model reliably achieves $H_{\text{stable}} \approx 4.09 \pm 0.19$ bits (10 seeds, reproduced by \texttt{experiments/verify\_table1\_preprint.py}), and component ablation identifies the doubt variable as the primary anti-synchronization mechanism: freezing $u$ raises pairwise synchrony from $0.007$ to $0.658$, a complete separation of the two conditions (Cohen's $d \approx 9$, 30-seed replication). Degree-normalized coupling extends diversity to sparse scale-free networks ($m \leq 3$ under the driven protocol; no power-law exponent $\gamma$ bridges the dead zone). Across twelve topology families the functional and dead-zone regimes separate empirically near $\lambda_2 \approx 2.31$, but degree-controlled experiments show this spectral separation to be correlational: the causal driver of the boundary is the coupling degree, through a mean-field sampling mechanism (harmonic degree $k_{\mathrm{harm}} \approx 6$), not algebraic connectivity.

The principal negative result is that for denser networks no degree-based reweighting preserves diversity, because averaging over more neighbours---not the multiplicity of spectral paths---collapses each node's target onto the global mean. Finite-size scaling of the Binder cumulant shows that the entropy loss across this boundary is a continuous crossover rather than a thermodynamic phase transition, and a circuit-level SPICE implementation reproduces the same functional/dead-zone dichotomy under analog fabrication tolerances. Unlike adaptive networks that tune coupling \emph{magnitude}, the mechanism here adapts coupling \emph{sign}, locally and per node; the failure of local corrections at high degree motivates future work on structural sparsification and node-specific drive rather than smarter weight normalization. More broadly, the boundary suggests a design constraint for neuromorphic and distributed systems: connectivity is not a free resource---past a critical coupling degree, additional links erase the very state heterogeneity that local mechanisms are meant to preserve, so architectures that must sustain functional diversity should bound the number of couplings per node rather than maximize it.

% ============================================================================
% 8. LIMITATIONS & TRANSPARENCY
% ============================================================================
\section{Limitations and Transparency}
\label{sec:limitations}

\subsection{Scientific Limitations}
\begin{itemize}
\item \textbf{Numerical integration:} All results use first-order Euler integration with $\Delta t = 0.05$. The spectral radius of the linearized Euler step reaches $\max|1 + \lambda_{\max}\,\Delta t| \approx 1.04 > 1$, which is formally outside the von~Neumann stability bound for the linearized system. In practice, the nonlinear restoring forces of the FitzHugh--Nagumo nullclines bound trajectories, preventing divergence. This tension between linear instability and nonlinear stability is a known feature of FHN-type systems. While Runge--Kutta 4 improves precision, Euler is sufficient for stability in our parameter regime ($\Delta t \le 0.05$), as validated by direct comparison ($\max \Delta H_{\text{cog}} < 0.006$).
\item \textbf{Heretic threshold:} The $\eta = 0.15$ value is empirically effective on 2D lattices ($k = 4$) but has not been validated as optimal across arbitrary topologies.
\item \textbf{Entropy metric and discretization:} Two entropy metrics are used in this work and must not be conflated. (1)~$H_{\text{stable}}$ (primary metric throughout) is a 100-bin continuous spectral entropy computed on the actual voltage distribution, recalibrated from earlier discrete estimates. (2)~$H_{\text{cog}}$ is a legacy 5-state discrete entropy using physiological bins $\{v < -1.2,\; {-1.2}\text{--}{-0.4},\; {-0.4}\text{--}{0.4},\; {0.4}\text{--}{1.2},\; v > 1.2\}$ inherited from the SPICE voltage scale. Python simulations at default parameters produce voltages predominantly in $[-3.2, -1.3]$, placing all nodes in bin~1 and yielding $H_{\text{cog}} \approx 0$ regardless of actual diversity. This is a calibration artefact, not a diversity collapse: pairwise synchrony ($\text{sync}_{\text{FULL}} = 0.007$ vs.\ $\text{sync}_{\text{FROZEN}} = 0.658$, a complete separation of the two distributions) and $H_{\text{stable}}$ confirm that diversity is real and substantial. $H_{\text{cog}}$ is retained only for compatibility with SPICE benchmarks, where voltage scales match the bin design.
\item \textbf{Doubt time-scale near the boundary:} The doubt recovery time constant $\tau_u = 10.0$ (default) was held fixed across all experiments. Near the connectivity boundary, the system's ability to sustain functional diversity may depend sensitively on $\tau_u$: a shorter $\tau_u$ could shift the effective threshold by allowing faster polarity reassignment, while a longer $\tau_u$ may produce critical-slowing-down signatures analogous to those observed near bifurcations in excitable media. The dependence of the critical coupling degree on $\tau_u$ has not been characterized and remains an open question.
\item \textbf{Adaptive coupling protocol --- unreconciled discrepancy:} Table~\ref{tab:adaptive_D_sweep} reports $H_{\text{cont}} \to 0.00$ for $D(u) = 0.50 \cdot u$ at $m \geq 5$ ($T = 10{,}000$ steps), whereas Table~\ref{tab:lz_regime} reports $\langle H_{\mathrm{stable}}\rangle = 4.23$ at $m = 6$ for the same protocol, topology and normalization ($T = 2000$ steps). The two tables thus support opposite readings---``catastrophic collapse by doubt saturation'' and ``structured dead zone''. The most likely difference is the integration length, which neither table presents as a factor. \emph{We report the discrepancy rather than resolve it}: the script that generated Table~\ref{tab:adaptive_D_sweep} is not in the repository, so a re-run would be a reconstruction rather than a verification, and could not distinguish an error in the table from a difference in our reconstruction. Both tables are secondary to the paper's claims, which rest on the ablation (Section~\ref{sec:ablation}) and the degree-controlled experiments (Section~\ref{sec:lambda2}).
\item \textbf{Scale:} The model has been tested up to $N = 2500$. Behavior at $N \gg 10{,}000$ is not validated.
\item \textbf{Connectivity regime boundary:} The boundary at $m \approx 6$ (endogenous, Cold Start Protocol) is empirically observed on $N = 100$ BA networks; it is governed by the coupling degree (Section~\ref{sec:lambda2}), and the exact critical harmonic degree ($k_{\mathrm{harm}} \approx 6$) depends on the regime metric ($H_{\mathrm{cog}}$) and may shift with parameters or at larger $N$.
\item \textbf{Replication count:} Primary results (Table~\ref{tab:scaling}, component ablation, scale-free topology sweep) and secondary validation experiments (spatial mutual information, stochastic resonance, $\gamma$ sweep) all use $n = 10$ seeds. The regime classification of Section~\ref{sec:lambda2} uses 3 seeds per topology family (36 rows total); as noted there, these labels are assigned per topology type rather than measured per seed, so the row count overstates the number of independent decisions. The causal claim rests on the degree-controlled experiments, not on this classification.
\end{itemize}

\subsection{Methodological Transparency}
This work was developed through a multi-agent collaborative research methodology involving multiple AI systems (Claude, GPT, Gemini) under continuous human guidance. All scientific claims, parameter choices, and publication decisions were made by the human author. AI contributions were advisory and generative. Complete session transcripts and intermediate versions are maintained in the public repository. The model underwent systematic adversarial testing via automated analysis of failure modes (December 2025--April 2026), leading to the corrections of entropy measurement (transient vs.\ sustained values) and the development of degree-normalized coupling.

% ============================================================================
% REFERENCES
% ============================================================================
\begin{thebibliography}{99}

% Consensus & Synchronization
\bibitem{olfati2007} Olfati-Saber, R., Fax, J.A., Murray, R.M., \emph{Consensus and cooperation in networked multi-agent systems}, Proceedings of the IEEE, 95(1):215-233, 2007.
\bibitem{strogatz2000} Strogatz, S.H., \emph{From Kuramoto to Crawford: Exploring the onset of synchronization in populations of coupled oscillators}, Physica D, 143(1-4):1-20, 2000.
\bibitem{acebrn2005} Acebr\'on, J.A., et al., \emph{The Kuramoto model: A simple paradigm for synchronization phenomena}, Reviews of Modern Physics, 77(1):137, 2005.

% FitzHugh-Nagumo
\bibitem{fitzhugh1961} FitzHugh, R., \emph{Impulses and physiological states in theoretical models of nerve membrane}, Biophysical Journal, 1(6):445-466, 1961.
\bibitem{nagumo1962} Nagumo, J., Arimoto, S., Yoshizawa, S., \emph{An active pulse transmission line simulating nerve axon}, Proceedings of the IRE, 50(10):2061-2070, 1962.
\bibitem{izhikevich2007} Izhikevich, E.M., \emph{Dynamical Systems in Neuroscience: The Geometry of Excitability and Bursting}, MIT Press, 2007.

% Network Science
\bibitem{barabasi1999} Barab\'asi, A.-L., Albert, R., \emph{Emergence of scaling in random networks}, Science, 286(5439):509-512, 1999.
\bibitem{watts1998} Watts, D.J., Strogatz, S.H., \emph{Collective dynamics of `small-world' networks}, Nature, 393(6684):440-442, 1998.

% Chimera States
\bibitem{abrams2004} Abrams, D.M., Strogatz, S.H., \emph{Chimera states for coupled oscillators}, Physical Review Letters, 93(17):174102, 2004.
\bibitem{omelchenko2013} Omelchenko, I., Omel'chenko, O.E., H\"ovel, P., Sch\"oll, E., \emph{When nonlocal coupling between oscillators becomes stronger: Patched synchrony or multichimera states}, Physical Review Letters, 110(22):224101, 2013.
\bibitem{omelchenko2015} Omelchenko, I., Provata, A., Hizanidis, J., Sch\"oll, E., H\"ovel, P., \emph{Robustness of chimera states for coupled FitzHugh--Nagumo oscillators}, Physical Review E, 91(2):022917, 2015.
\bibitem{semenova2016} Semenova, N., Zakharova, A., Anishchenko, V., Sch\"oll, E., \emph{Coherence-resonance chimeras in a network of excitable elements}, Physical Review Letters, 117(1):014102, 2016.
\bibitem{majhi2019} Majhi, S., Bera, B.K., Ghosh, D., Perc, M., \emph{Chimera states in neuronal networks: A review}, Physics of Life Reviews, 28:100--121, 2019.
\bibitem{provata2026} Provata, A., Boulougouris, G.C., Hizanidis, J., \emph{Adaptive transitions in FitzHugh--Nagumo networks with Hebb--Oja coupling rules}, arXiv:2602.18198, 2026.
\bibitem{berner2023} Berner, R., Gross, T., Kuehn, C., Kurths, J., Yanchuk, S., \emph{Adaptive dynamical networks}, Physics Reports, 1031:1--59, 2023.
\bibitem{zanette2004} Zanette, D.H., Mikhailov, A.S., \emph{Dynamical systems with time-dependent coupling: Clustering and critical behaviour}, Physica D, 194(3--4):203--218, 2004.
\bibitem{hong2011} Hong, H., Strogatz, S.H., \emph{Kuramoto model of coupled oscillators with positive and negative coupling parameters: An example of conformist and contrarian oscillators}, Physical Review Letters, 106(5):054102, 2011.

% Opinion Dynamics
\bibitem{castellano2009} Castellano, C., Fortunato, S., Loreto, V., \emph{Statistical physics of social dynamics}, Reviews of Modern Physics, 81(2):591, 2009.
\bibitem{hegselmann2002} Hegselmann, R., Krause, U., \emph{Opinion dynamics and bounded confidence models, analysis, and simulation}, Journal of Artificial Societies and Social Simulation, 5(3), 2002.

% Multi-Agent Systems
\bibitem{jadbabaie2003} Jadbabaie, A., Lin, J., Morse, A.S., \emph{Coordination of groups of mobile autonomous agents using nearest neighbor rules}, IEEE Transactions on Automatic Control, 48(6):988-1001, 2003.

% Nonlinear Dynamics
\bibitem{pikovsky2001} Pikovsky, A., Rosenblum, M., Kurths, J., \emph{Synchronization: A Universal Concept in Nonlinear Sciences}, Cambridge University Press, 2001.

% Frustrated Synchronization & Spin Glasses
\bibitem{toulouse1977} Toulouse, G., \emph{Theory of the frustration effect in spin glasses: I}, Communications on Physics, 2:115--119, 1977.
\bibitem{odor2024} \'Odor, G., Deng, S., Kelling, J., \emph{Frustrated Synchronization of the Kuramoto Model on Complex Networks}, Entropy, 26(12):1074, 2024.

% Random-Field Ising Model
\bibitem{nattermann1998} Nattermann, T., \emph{Theory of the Random Field Ising Model}, in \emph{Spin Glasses and Random Fields}, ed.\ A.P.\ Young, World Scientific, pp.\ 277--298, 1998.

% Adaptive Networks
\bibitem{gross2008} Gross, T., Blasius, B., \emph{Adaptive coevolutionary networks: a review}, Journal of the Royal Society Interface, 5(20):259-271, 2008.

% Statistics
\bibitem{albert1984} Albert, A., Anderson, J.A., \emph{On the existence of maximum likelihood estimates in logistic regression models}, Biometrika, 71(1):1--10, 1984.
\bibitem{binder1981finite} Binder, K., \emph{Finite size scaling analysis of ising model block distribution functions}, Zeitschrift f\"ur Physik B Condensed Matter, 43(2):119-140, 1981.

% Computational Neuroscience / Active Inference
\bibitem{friston2010} Friston, K., \emph{The free-energy principle: a unified brain theory?}, Nature Reviews Neuroscience, 11(2):127--138, 2010.

% AI Methodology
\bibitem{bommasani2021} Bommasani, R., et al., \emph{On the Opportunities and Risks of Foundation Models}, arXiv:2108.07258, 2021.
\bibitem{park2023} Park, J.S., et al., \emph{Generative Agents: Interactive Simulacra of Human Behavior}, Proceedings of the 36th Annual ACM Symposium on UIST, 2023.

\end{thebibliography}

% ============================================================================
% APPENDICES
% ============================================================================
\appendix

\section{Full Algorithm}
\label{app:algorithm}

\begin{algorithm}[ht]
\caption{Time-step update for the doubt-modulated FHN network.}
\begin{algorithmic}[1]
\FOR{each time-step $t$}
    \STATE \COMMENT{Optional: Doubt-driven rewiring (explicit graphs only)}
    \FOR{each unit $i$ with $u_i > u_{\text{rewire}}$ and cooldown expired}
        \STATE $j^* \leftarrow \arg\min_{k \in \mathcal{N}(i)} |v_i - v_k|$
        \STATE $k \leftarrow \text{RandomNonNeighbor}(i)$
        \STATE Disconnect $i \leftrightarrow j^*$; Connect $i \leftrightarrow k$
    \ENDFOR
    \FOR{each unit $i$}
        \STATE $\sigma_{\text{local},i} \leftarrow |\bar{v}_{\mathcal{N}(i)} - v_i|$
        \STATE $\eta_v \leftarrow \mathcal{N}(0, \sigma_{\text{noise}}^2)$
        \STATE $w_i^{\text{coup}} \leftarrow \tanh(\pi(0.5 - u_i)) + \delta$
        \STATE $I_{\text{coup}} \leftarrow D_{\text{eff}}(i) \cdot w_i^{\text{coup}} \cdot (\bar{v}_{\mathcal{N}(i)} - v_i)$
        \STATE $I_{\text{ext}} \leftarrow s_i \cdot I_{\text{stimulus}} + I_{\text{coup}}$ \COMMENT{$s_i = -1$ for heretics}
        \STATE $v_i \leftarrow v_i + (v_i - v_i^3/5 - w_i + I_{\text{ext}} - \alpha_{\text{self}} \tanh(v_i) + \eta_v) \cdot \Delta t$
        \STATE $w_i \leftarrow w_i + (\varepsilon(v_i + a - bw_i) + dw_{\text{plast},i}) \cdot \Delta t$
        \STATE $\varepsilon_{u,\text{eff}} \leftarrow \varepsilon_u \cdot \min(1 + \alpha_s \sigma_{\text{local},i},\; C_{\text{cap}})$
        \STATE $u_i \leftarrow \text{clamp}\!\bigl(u_i + \frac{\varepsilon_{u,\text{eff}}}{\tau_u}(k_u \sigma_{\text{local},i} + \sigma_{\text{base}} - u_i) \cdot \Delta t,\; 0,\; 1\bigr)$
    \ENDFOR
    \STATE Compute $H(t)$ from state distribution (Table~\ref{tab:states})
\ENDFOR
\end{algorithmic}
\end{algorithm}

\section{Legacy Cognitive State Bins}
\label{app:legacy_bins}

The 5-state discrete entropy $H_{\text{cog}}$ uses the physiological voltage bins of Table~\ref{tab:states}, inherited from the SPICE voltage scale. It is retained for circuit-level comparisons only (see Limitations, Section~\ref{sec:limitations}).

\begin{table}[ht]
\centering
\caption{Cognitive state classification by activation potential $v$.}
\label{tab:states}
\begin{tabular}{lcc}
\toprule
\textbf{State} & \textbf{Range} & \textbf{Interpretation} \\
\midrule
1 (Strong negative) & $v < -1.2$ & Strong opposition \\
2 (Weak negative)   & $-1.2 \leq v < -0.4$ & Mild opposition \\
3 (Neutral)         & $-0.4 \leq v \leq 0.4$ & Uncertain / Deliberative \\
4 (Weak positive)   & $0.4 < v \leq 1.2$ & Mild agreement \\
5 (Strong positive) & $v > 1.2$ & Strong agreement \\
\bottomrule
\end{tabular}
\end{table}

\section{Reference Parameters}
\label{app:parameters}

\begin{table}[ht]
\centering
\caption{Default parameter values. All parameters are configurable via YAML.}
\begin{tabular}{llrl}
\toprule
\textbf{Group} & \textbf{Parameter} & \textbf{Value} & \textbf{Description} \\
\midrule
Dynamics & $a$ & 0.7 & FHN offset \\
         & $b$ & 0.8 & Recovery rate \\
         & $\varepsilon$ & 0.08 & Time-scale separation \\
         & $\alpha_{\text{self}}$ & 0.15 & Self-inhibition gain \\
         & $\Delta t$ & 0.05 & Integration step \\
\midrule
Coupling & $D$ & 0.15 & Base coupling strength \\
         & $\delta$ & 0.01 & Sigmoid offset \\
         & $\eta$ & 0.15 & Heretic fraction \\
\midrule
Doubt    & $\varepsilon_u$ & 0.02 & Doubt learning rate \\
         & $\tau_u$ & 10.0 & Doubt time constant \\
         & $k_u$ & 1.0 & Local disagreement gain \\
         & $\sigma_{\text{baseline}}$ & 0.05 & Minimum doubt \\
         & $u_{\text{clamp}}$ & $[0, 1]$ & Doubt range \\
\midrule
Meta-doubt & $\alpha_s$ & 2.0 & Surprise gain \\
           & $C_{\text{cap}}$ & 5.0 & Acceleration cap \\
\midrule
Noise    & $\sigma_{\text{noise}}$ & 0.05 & State noise std \\
\midrule
Rewiring & $u_{\text{rewire}}$ & 0.8 & Rewire threshold \\
         & Cooldown & 50 & Steps between rewires \\
\bottomrule
\end{tabular}
\end{table}

\end{document}
